
\Data\ID{1003187301}
\Updated{2020-12-29 08:12:21}
\Level{A}
\SubArea{Absolute value equations and inequalities}
\LANGen{
\Question{
Which equation has the solutions \( -1 \) and \( 9 \)?}
{\( |2x-8|=10 \)}{\( |2x-10|=8 \)}{\( |2x+8|=10 \)}{\( |2x+10|=8 \)}{}{}}
\LANGpl{
\Question{
Dla, którego równania rozwiązaniem są  \( -1 \) i \( 9 \)?}
{\( |2x-8|=10 \)}{\( |2x-10|=8 \)}{\( |2x+8|=10 \)}{\( |2x+10|=8 \)}{}{}}
\LANGcs{
\Question{
Která rovnice má řešení \( -1 \) a \( 9 \)?}
{\( |2x-8|=10 \)}{\( |2x-10|=8 \)}{\( |2x+8|=10 \)}{\( |2x+10|=8 \)}{}{}}
\LANGsk{
\Question{
Ktorá rovnica má riešenie \( -1 \) a \( 9 \)?}
{\( |2x-8|=10 \)}{\( |2x-10|=8 \)}{\( |2x+8|=10 \)}{\( |2x+10|=8 \)}{}{}}
\LANGes{
\Question{
¿Cuál de las ecuaciones siguientes tiene como soluciones \( -1 \) y \( 9 \)?}
{\( |2x-8|=10 \)}{\( |2x-10|=8 \)}{\( |2x+8|=10 \)}{\( |2x+10|=8 \)}{}{}}
\End

\Data\ID{1003187306}
\Updated{2020-08-23 01:08:58}
\Level{A}
\SubArea{Absolute value equations and inequalities}
\LANGen{
\Question{
Let \( |2x+6|-4=0 \). Choose the correct statement.}
{The equation has two solutions.}{Any real value of \( x \) is a solution.}{The equation has just one solution.}{The equation has no solution.}{}{}}
\LANGpl{
\Question{
Niech \( |2x+6|-4=0 \). Wybierz prawdziwe stwierdzenie.}
{Równanie ma dwa rozwiązania.}{Każda rzeczywista liczba \( x \) jest rozwiązaniam.}{Równanie ma tylko jedno rozwiązanie.}{Równanie nie ma rozwiązania.}{}{}}
\LANGcs{
\Question{
Nechť \( |2x+6|-4=0 \). Vyberte správný výrok.}
{Rovnice má dvě řešení.}{Jakákoliv reálná hodnota \( x \) je řešením.}{Rovnice má právě jedno řešení.}{Rovnice nemá řešení.}{}{}}
\LANGsk{
\Question{
Nech \( |2x+6|-4=0 \). Vyberte správny výrok.}
{Rovnica má dve riešenia.}{Akákoľvek reálna hodnota \( x \) je riešením.}{Rovnica má práve jedno riešenie.}{Rovnica nemá riešenie.}{}{}}
\LANGes{
\Question{
Dada la ecuación \( |2x+6|-4=0 \). Elige el enunciado correcto.}
{La ecuación tiene dos soluciones.}{Cualquier valor real de \( x \) es la solución.}{La ecuación tiene solo una solución.}{La ecuación no tiene solución.}{}{}}
\End

\Data\ID{1003187307}
\Updated{2020-08-23 07:08:33}
\Level{A}
\SubArea{Absolute value equations and inequalities}
\LANGen{
\Question{
The solutions of the equation \( |14-4x|=4 \) are:}
{the numbers differing by \( 2 \).}{the numbers differing by \( 1 \).}{integers.}{opposite numbers.}{}{}}
\LANGpl{
\Question{
Rozwiązaniem równania \( |14-4x|=4 \) są:}
{liczby różniące się o \( 2 \).}{liczby różniące się o \( 1 \).}{liczby całkowite.}{odwrotne liczby.}{}{}}
\LANGcs{
\Question{
Kořeny rovnice \( |14-4x|=4 \) jsou:}
{čísla, jejichž rozdíl je roven \( 2 \).}{čísla, jejichž rozdíl je roven \( 1 \).}{celá čísla.}{opačná čísla.}{}{}}
\LANGsk{
\Question{
Korene rovnice \( |14-4x|=4 \) sú:}
{čísla, ktorých rozdiel sa rovná  \( 2 \).}{čísla, ktorých rozdiel sa rovná \( 1 \).}{celé čísla.}{opačné čísla.}{}{}}
\LANGes{
\Question{
Las soluciones de la ecuación \( |14-4x|=4 \) son:}
{números cuyo resto es igual a \( 2 \).}{números cuyo resto es igual a \( 1 \).}{números enteros.}{números opuestos.}{}{}}
\End

\Data\ID{1003187308}
\Updated{2022-04-23 05:04:37}
\Level{A}
\SubArea{Absolute value equations and inequalities}
\LANGen{
\Question{
Choose the equation which has only one solution.}
{\( 4+|2x-6|=4 \)}{\(  2-|x-3|=1 \)}{\( |x-3|+2=-1 \)}{\( 2-|x-3|=-2 \)}{}{}}
\LANGpl{
\Question{
Wybierz równanie, które ma tylko jedno rozwiązanie.}
{\( 4+|2x-6|=4 \)}{\(  2-|x-3|=1 \)}{\( |x-3|+2=-1 \)}{\( 2-|x-3|=-2 \)}{}{}}
\LANGcs{
\Question{
Vyberte rovnici, která má jen jedno řešení.}
{\( 4+|2x-6|=4 \)}{\(  2-|x-3|=1 \)}{\( |x-3|+2=-1 \)}{\( 2-|x-3|=-2 \)}{}{}}
\LANGsk{
\Question{
Vyberte rovnicu, ktorá má len jedno riešenie.}
{\( 4+|2x-6|=4 \)}{\(  2-|x-3|=1 \)}{\( |x-3|+2=-1 \)}{\( 2-|x-3|=-2 \)}{}{}}
\LANGes{
\Question{
Elige la ecuación que tiene solo una solución.}
{\( 4+|2x-6|=4 \)}{\(  2-|x-3|=1 \)}{\( |x-3|+2=-1 \)}{\( 2-|x-3|=-2 \)}{}{}}
\End

\Data\ID{1003187312}
\Updated{2020-12-29 08:12:24}
\Level{A}
\SubArea{Absolute value equations and inequalities}
\LANGen{
\Question{
The solution set of an inequality in interval notation is \( (-\infty;-12]\cup[12;\infty) \).  Find this inequality.}
{\( |x| \geq 12 \)}{\( |x|\leq 12 \)}{\( |x| > 12 \)}{\( |x| < 12 \)}{}{}}
\LANGpl{
\Question{
Zbiór rozwiązań nierówności mieści się w przedziale \( (-\infty;-12\rangle\cup\langle12;\infty) \).  Wyznacz tę nierówność.}
{\( |x| \geq 12 \)}{\( |x|\leq 12 \)}{\( |x| > 12 \)}{\( |x| < 12 \)}{}{}}
\LANGcs{
\Question{
Množina řešení nerovnice v zápisu intervalu je \( (-\infty;-12\rangle\cup\langle12;\infty) \).  Určete nerovnici.}
{\( |x| \geq 12 \)}{\( |x|\leq 12 \)}{\( |x| > 12 \)}{\( |x| < 12 \)}{}{}}
\LANGsk{
\Question{
Množina riešení nerovnice v zápise intervalu je \( (-\infty;-12\rangle\cup\langle12;\infty) \).  Určte nerovnicu.}
{\( |x| \geq 12 \)}{\( |x|\leq 12 \)}{\( |x| > 12 \)}{\( |x| < 12 \)}{}{}}
\LANGes{
\Question{
El conjunto de soluciones de una inecuación viene dado por el intervalo \( (-\infty;-12]\cup[12;\infty) \). Encuentra dicha inecuación.}
{\( |x| \geq 12 \)}{\( |x|\leq 12 \)}{\( |x| > 12 \)}{\( |x| < 12 \)}{}{}}
\End

\Data\ID{1003187406}
\Updated{2022-04-23 05:04:48}
\Level{A}
\SubArea{Absolute value equations and inequalities}
\LANGen{
\Question{
How many natural numbers belong to the solution set of the inequality \( |4x-10| \leq 6 \)?}
{\( 4 \)}{\( 2 \)}{\( 3 \)}{\( 0 \)}{}{}}
\LANGpl{
\Question{
Ile liczb naturalnych należy do zbioru rozwiązań nierówności \( |4x-10| \leq 6 \)?}
{\( 4 \)}{\( 2 \)}{\( 3 \)}{\( 0 \)}{}{}}
\LANGcs{
\Question{
Kolik přirozených čísel patří do množiny řešení nerovnice \( |4x-10| \leq 6 \)?}
{\( 4 \)}{\( 2 \)}{\( 3 \)}{\( 0 \)}{}{}}
\LANGsk{
\Question{
Koľko prirodzených čísel patrí množine riešení nerovnice \( |4x-10| \leq 6 \)?}
{\( 4 \)}{\( 2 \)}{\( 3 \)}{\( 0 \)}{}{}}
\LANGes{
\Question{
¿Cuántos números naturales pertenecen al conjunto de soluciones de la inecuación \( |4x-10| \leq 6 \)?}
{\( 4 \)}{\( 2 \)}{\( 3 \)}{\( 0 \)}{}{}}
\End

\Data\ID{1003187407}
\Updated{2022-04-23 05:04:48}
\Level{A}
\SubArea{Absolute value equations and inequalities}
\LANGen{
\Question{
Find the solution set of the inequality \( |12-6x| < 2 \).}
{\( \left( \frac53;\frac73 \right) \)}{\( \left( -\frac73;-\frac53 \right) \)}{\( \left( -\frac73;\frac53 \right) \)}{\( \left( -\frac53;\frac73 \right) \)}{}{}}
\LANGpl{
\Question{
Wyznacz zbiór rozwiązań nierówności \( |12-6x| < 2 \).}
{\( \left( \frac53;\frac73 \right) \)}{\( \left( -\frac73;-\frac53 \right) \)}{\( \left( -\frac73;\frac53 \right) \)}{\( \left( -\frac53;\frac73 \right) \)}{}{}}
\LANGcs{
\Question{
Určete množinu řešení nerovnice \( |12-6x| < 2 \).}
{\( \left( \frac53;\frac73 \right) \)}{\( \left( -\frac73;-\frac53 \right) \)}{\( \left( -\frac73;\frac53 \right) \)}{\( \left( -\frac53;\frac73 \right) \)}{}{}}
\LANGsk{
\Question{
Určte množinu riešení nerovnice \( |12-6x| < 2 \).}
{\( \left( \frac53;\frac73 \right) \)}{\( \left( -\frac73;-\frac53 \right) \)}{\( \left( -\frac73;\frac53 \right) \)}{\( \left( -\frac53;\frac73 \right) \)}{}{}}
\LANGes{
\Question{
Halla el conjunto de soluciones de la inecuación \( |12-6x| < 2 \).}
{\( \left( \frac53;\frac73 \right) \)}{\( \left( -\frac73;-\frac53 \right) \)}{\( \left( -\frac73;\frac53 \right) \)}{\( \left( -\frac53;\frac73 \right) \)}{}{}}
\End

\Data\ID{1103187310}
\Updated{2020-12-29 08:12:01}
\Level{A}
\SubArea{Absolute value equations and inequalities}
\IMG{1103187310.pdf}
\LANGen{
\Question{
The solution set of an inequality is graphed on the number line. Determine this inequality.}
{\( |x+2| \leq 3 \)}{\( |x-2| \leq 3 \)}{\( |x-3| \leq 2 \)}{\( |x+3| \leq 2 \)}{}{}}
\LANGpl{
\Question{
Zbiór rozwiązań nierówności przedstawiony jest za pomocą linii liczbowej. Wyznacz tę nierówność.}
{\( |x+2| \leq 3 \)}{\( |x-2| \leq 3 \)}{\( |x-3| \leq 2 \)}{\( |x+3| \leq 2 \)}{}{}}
\LANGcs{
\Question{
Množina řešení nerovnice je znázorněno na číselné ose. Určete tuto nerovnici.}
{\( |x+2| \leq 3 \)}{\( |x-2| \leq 3 \)}{\( |x-3| \leq 2 \)}{\( |x+3| \leq 2 \)}{}{}}
\LANGsk{
\Question{
Množina riešení nerovnice je znázornená na číselnej osi. Určte túto nerovnicu.}
{\( |x+2| \leq 3 \)}{\( |x-2| \leq 3 \)}{\( |x-3| \leq 2 \)}{\( |x+3| \leq 2 \)}{}{}}
\LANGes{
\Question{
En la imagen se representa el conjunto de soluciones de una inecuación. Define esta inecuación.}
{\( |x+2| \leq 3 \)}{\( |x-2| \leq 3 \)}{\( |x-3| \leq 2 \)}{\( |x+3| \leq 2 \)}{}{}}
\End

\Data\ID{1103187311}
\Updated{2020-08-23 01:08:27}
\Level{A}
\SubArea{Absolute value equations and inequalities}
\LANGen{
\Question{
Which graph shows the solution set of the inequality \( |x-3| \geq 1 \)?}
{}{}{}{}{}{}}
\LANGpl{
\Question{
Który wykres przedstawia zbiór rozwiązań nierówności  \( |x-3| \geq 1 \)?}
{}{}{}{}{}{}}
\LANGcs{
\Question{
Který graf znázorňuje množinu řešení nerovnice \( |x-3| \geq 1 \)?}
{}{}{}{}{}{}}
\LANGsk{
\Question{
Ktorý graf znázorňuje množinu riešení nerovnice \( |x-3| \geq 1 \)?}
{}{}{}{}{}{}}
\LANGes{
\Question{
¿Cuál de las gráficas representa el conjunto de soluciones de la inecuación \( |x-3| \geq 1 \)?}
{}{}{}{}{}{}}

\IMGanswer{1103187311a.pdf}
\IMGanswer{1103187311b.pdf}
\IMGanswer{1103187311c.pdf}
\IMGanswer{1103187311d.pdf}\End

\Data\ID{1103187404}
\Updated{2022-04-23 05:04:48}
\Level{A}
\SubArea{Absolute value equations and inequalities}
\IMG{1103187404.pdf}
\LANGen{
\Question{
The solution set of an inequality is graphed on the number line. Determine this inequality.}
{\( |4-x| > 41 \)}{\( |x-3| < 42 \)}{\( |x-2| > 42 \)}{\( |1-x| > 43 \)}{}{}}
\LANGpl{
\Question{
Zbiór rozwiązań nierówności przedstawiono za pomocą linii liczbowej. Wyznacz tę nierówność.}
{\( |4-x| > 41 \)}{\( |x-3| < 42 \)}{\( |x-2| > 42 \)}{\( |1-x| > 43 \)}{}{}}
\LANGcs{
\Question{
Řešení množiny nerovnice je znázorněno na číselné ose. Určete tuto nerovnici.}
{\( |4-x| > 41 \)}{\( |x-3| < 42 \)}{\( |x-2| > 42 \)}{\( |1-x| > 43 \)}{}{}}
\LANGsk{
\Question{
Množina riešení nerovnice je znázornená na číselnej osi. Určte túto nerovnicu.}
{\( |4-x| > 41 \)}{\( |x-3| < 42 \)}{\( |x-2| > 42 \)}{\( |1-x| > 43 \)}{}{}}
\LANGes{
\Question{
En la imagen se representa el conjunto de soluciones de una inecuación. Define esta inecuación.}
{\( |4-x| > 41 \)}{\( |x-3| < 42 \)}{\( |x-2| > 42 \)}{\( |1-x| > 43 \)}{}{}}
\End

\Data\ID{1103187408}
\Updated{2020-08-23 01:08:45}
\Level{A}
\SubArea{Absolute value equations and inequalities}
\LANGen{
\Question{
Identify the graph which shows the solution set of the inequality \( |2x-12| > 6 \).}
{}{}{}{}{}{}}
\LANGpl{
\Question{
Wybierz wykres, który przedstawia zbiór rozwiązań nierówności \( |2x-12| > 6 \).}
{}{}{}{}{}{}}
\LANGcs{
\Question{
Určete graf, který znázorňuje množinu řešení nerovnice \( |2x-12| > 6 \).}
{}{}{}{}{}{}}
\LANGsk{
\Question{
Určte graf, ktorý znázorňuje množinu riešení nerovnice \( |2x-12| > 6 \).}
{}{}{}{}{}{}}
\LANGes{
\Question{
Identifica la gráfica que representa el conjunto de soluciones de la inecuación \( |2x-12| > 6 \).}
{}{}{}{}{}{}}

\IMGanswer{1103187408a.pdf}
\IMGanswer{1103187408b.pdf}
\IMGanswer{1103187408c.pdf}
\IMGanswer{1103187408d.pdf}\End

\Data\ID{1103187409}
\Updated{2022-04-23 05:04:48}
\Level{A}
\SubArea{Absolute value equations and inequalities}
\IMG{1103187409.pdf}
\LANGen{
\Question{
The graph shows the set of all real numbers that satisfy the inequality
\( |3x-12| \leq 15 \). Determine \( k \).}
{\( k = 9 \)}{\( k = 4 \)}{\( k = 5 \)}{\( k = 2 \)}{}{}}
\LANGpl{
\Question{
Wykres przedstawia zbiór wszystkich rzeczywistych liczb, które spełniają nierówność
\( |3x-12| \leq 15 \). Wyznacz \( k \).}
{\( k = 9 \)}{\( k = 4 \)}{\( k = 5 \)}{\( k = 2 \)}{}{}}
\LANGcs{
\Question{
Graf znázorňuje množinu všech reálných čísel platných pro nerovnici
\( |3x-12| \leq 15 \). Určete \( k \).}
{\( k = 9 \)}{\( k = 4 \)}{\( k = 5 \)}{\( k = 2 \)}{}{}}
\LANGsk{
\Question{
Graf znázorňuje množinu všetkých reálnych čísel platných pre nerovnicu 
\( |3x-12| \leq 15 \). Určte \( k \).}
{\( k = 9 \)}{\( k = 4 \)}{\( k = 5 \)}{\( k = 2 \)}{}{}}
\LANGes{
\Question{
La gráfica representa el conjunto de todos los números reales que cumplen la inecuación 
\( |3x-12| \leq 15 \). Determina \( k \).}
{\( k = 9 \)}{\( k = 4 \)}{\( k = 5 \)}{\( k = 2 \)}{}{}}
\End

\Data\ID{2000001101}
\Updated{2021-12-01 03:12:46}
\Level{A}
\SubArea{Absolute value equations and inequalities}
\LANGen{
\Question{
Find the solution set of the equation \(|x + 1| = 0\).}
{\(\{-1\}\)}{\(\emptyset\)}{\(\{-1;1\}\)}{\(\{1\}\)}{}{}}
\LANGpl{
\Question{
Wyznacz zbiór rozwiązań równania \(|x + 1| = 0\).}
{\(\{-1\}\)}{\(\emptyset\)}{\(\{-1;1\}\)}{\(\{1\}\)}{}{}}
\LANGcs{
\Question{
Určete množinu řešení rovnice \(|x + 1| = 0\).}
{\(\{-1\}\)}{\(\emptyset\)}{\(\{-1;1\}\)}{\(\{1\}\)}{}{}}
\LANGsk{
\Question{
Určte množinu riešení rovnice \(|x + 1| = 0\).}
{\(\{-1\}\)}{\(\emptyset\)}{\(\{-1;1\}\)}{\(\{1\}\)}{}{}}
\LANGes{
\Question{
Halla el conjunto de soluciones de la ecuación \(|x + 1| = 0\).}
{\(\{-1\}\)}{\(\emptyset\)}{\(\{-1;1\}\)}{\(\{1\}\)}{}{}}
\End

\Data\ID{2000001102}
\Updated{2021-12-01 03:12:41}
\Level{A}
\SubArea{Absolute value equations and inequalities}
\LANGen{
\Question{
Find the solution set of the equation \( |x-1|=12 \).}
{\(\{-11;13\}\)}{\(\{-13;11\}\)}{\(\{-11\}\)}{\(\{13\}\)}{}{}}
\LANGpl{
\Question{
Wyznacz zbiór rozwiązań równania \( |x-1|=12 \).}
{\(\{-11;13\}\)}{\(\{-13;11\}\)}{\(\{-11\}\)}{\(\{13\}\)}{}{}}
\LANGcs{
\Question{
Určete množinu řešení rovnice \( |x-1|=12 \).}
{\(\{-11;13\}\)}{\(\{-13;11\}\)}{\(\{-11\}\)}{\(\{13\}\)}{}{}}
\LANGsk{
\Question{
Určte množinu riešení rovnice \( |x-1|=12 \).}
{\(\{-11;13\}\)}{\(\{-13;11\}\)}{\(\{-11\}\)}{\(\{13\}\)}{}{}}
\LANGes{
\Question{
Halla el conjunto de soluciones de la ecuación \( |x-1|=12 \).}
{\(\{-11;13\}\)}{\(\{-13;11\}\)}{\(\{-11\}\)}{\(\{13\}\)}{}{}}
\End

\Data\ID{2000001103}
\Updated{2021-12-01 03:12:21}
\Level{A}
\SubArea{Absolute value equations and inequalities}
\LANGen{
\Question{
Choose the equation that has only positive solutions.}
{\( |x-5|=3\)}{\( |x-3|=5\)}{\( |x-2|=2\)}{\( |x+10|=8\)}{}{}}
\LANGpl{
\Question{
Wybierz równanie, które ma tylko jedno poprawne rozwiązanie.}
{\( |x-5|=3\)}{\( |x-3|=5\)}{\( |x-2|=2\)}{\( |x+10|=8\)}{}{}}
\LANGcs{
\Question{
Vyberte rovnici, pro kterou platí, že všechny její kořeny jsou kladná čísla.}
{\( |x-5|=3\)}{\( |x-3|=5\)}{\( |x-2|=2\)}{\( |x+10|=8\)}{}{}}
\LANGsk{
\Question{
Vyberte rovnicu, pre ktorú platí, že všetky jej korene sú kladné čísla.}
{\( |x-5|=3\)}{\( |x-3|=5\)}{\( |x-2|=2\)}{\( |x+10|=8\)}{}{}}
\LANGes{
\Question{
Elige la ecuación que tiene solo soluciones positivas.}
{\( |x-5|=3\)}{\( |x-3|=5\)}{\( |x-2|=2\)}{\( |x+10|=8\)}{}{}}
\End

\Data\ID{2000001104}
\Updated{2021-12-01 03:12:16}
\Level{A}
\SubArea{Absolute value equations and inequalities}
\LANGen{
\Question{
Choose the equation that has no solution.}
{\( |3x-5|=-1 \)}{\( |3x-5|=0\)}{\( |3x-5|=1 \)}{\( 3|x+5|=0 \)}{}{}}
\LANGpl{
\Question{
Wybierz równanie, które nie ma rozwiązania.}
{\( |3x-5|=-1 \)}{\( |3x-5|=0\)}{\( |3x-5|=1 \)}{\( 3|x+5|=0 \)}{}{}}
\LANGcs{
\Question{
Vyberte rovnici, která nemá řešení.}
{\( |3x-5|=-1 \)}{\( |3x-5|=0\)}{\( |3x-5|=1 \)}{\( 3|x+5|=0 \)}{}{}}
\LANGsk{
\Question{
Vyberte rovnicu, ktorá nemá riešenie.}
{\( |3x-5|=-1 \)}{\( |3x-5|=0\)}{\( |3x-5|=1 \)}{\( 3|x+5|=0 \)}{}{}}
\LANGes{
\Question{
Elige la ecuación que no tiene solución.}
{\( |3x-5|=-1 \)}{\( |3x-5|=0\)}{\( |3x-5|=1 \)}{\( 3|x+5|=0 \)}{}{}}
\End

\Data\ID{2000001105}
\Updated{2021-12-01 03:12:11}
\Level{A}
\SubArea{Absolute value equations and inequalities}
\LANGen{
\Question{
Choose the equation that has only one solution.}
{\(- |x-3|=0 \)}{\( |x-3|+1=0 \)}{\( |x-3|=3 \)}{\( 3+|x-3|=0 \)}{}{}}
\LANGpl{
\Question{
Wybierz równanie, które ma jedno rozwiązanie.}
{\(- |x-3|=0 \)}{\( |x-3|+1=0 \)}{\( |x-3|=3 \)}{\( 3+|x-3|=0 \)}{}{}}
\LANGcs{
\Question{
Vyberte rovnici, která má jen jedno řešení.}
{\(- |x-3|=0 \)}{\( |x-3|+1=0 \)}{\( |x-3|=3 \)}{\( 3+|x-3|=0 \)}{}{}}
\LANGsk{
\Question{
Vyberte rovnicu, ktorá má len jedno riešenie.}
{\(- |x-3|=0 \)}{\( |x-3|+1=0 \)}{\( |x-3|=3 \)}{\( 3+|x-3|=0 \)}{}{}}
\LANGes{
\Question{
Elige la ecuación que tiene solo una solución.}
{\(- |x-3|=0 \)}{\( |x-3|+1=0 \)}{\( |x-3|=3 \)}{\( 3+|x-3|=0 \)}{}{}}
\End

\Data\ID{2000002201}
\Updated{2021-12-01 03:12:26}
\Level{A}
\SubArea{Absolute value equations and inequalities}
\LANGen{
\Question{
Identify the solution set of the following inequality.
\[
|2x+5|>0
\]}
{\( \mathbb{R} \setminus \{-2.5\} \)}{\( \emptyset \)}{\( (-2.5; \infty) \)}{\( (-\infty;-2.5 ) \)}{}{}}
\LANGpl{
\Question{
Wyznacz zbiór rozwiązań następującej nierówności. 
\[
|2x+5|>0
\]}
{\( \mathbb{R} \setminus \{-2{,}5\} \)}{\( \emptyset \)}{\( (-2{,}5; \infty) \)}{\( (-\infty;-2{,}5 ) \)}{}{}}
\LANGcs{
\Question{
Vyberte množinu, která je řešením dané nerovnice.
\[
|2x+5|>0
\]}
{\( \mathbb{R} \setminus \{-2{,}5\} \)}{\( \emptyset \)}{\( (-2{,}5; \infty) \)}{\( (-\infty;-2{,}5 ) \)}{}{}}
\LANGsk{
\Question{
Vyberte množinu, ktorá je riešením danej  nerovnice.
\[
|2x+5|>0
\]}
{\( \mathbb{R} \setminus \{-2{,}5\} \)}{\( \emptyset \)}{\( (-2{,}5; \infty) \)}{\( (-\infty;-2{,}5 ) \)}{}{}}
\LANGes{
\Question{
Identifica el conjunto de soluciones de la siguiente inecuación. \[
|2x+5|>0
\]}
{\( \mathbb{R} \setminus \{-2.5\} \)}{\( \emptyset \)}{\( (-2.5; \infty) \)}{\( (-\infty;-2.5 ) \)}{}{}}
\End

\Data\ID{2000002202}
\Updated{2021-12-01 03:12:21}
\Level{A}
\SubArea{Absolute value equations and inequalities}
\LANGen{
\Question{
Identify the solution set of the following inequality.
\[
|x-25| \leq 0
\]}
{\( \{25\} \)}{\( \emptyset \)}{\( \mathbb{R} \setminus \{25\} \)}{\( (-\infty;-25) \)}{}{}}
\LANGpl{
\Question{
Wyznacz zbiór rozwiązań następującej nierówności. 
\[
|x-25| \leq 0
\]}
{\( \{25\} \)}{\( \emptyset \)}{\( \mathbb{R} \setminus \{25\} \)}{\( (-\infty;-25) \)}{}{}}
\LANGcs{
\Question{
Vyberte množinu, která je řešením dané nerovnice.
\[
|x-25| \leq 0
\]}
{\( \{25\} \)}{\( \emptyset \)}{\( \mathbb{R} \setminus \{25\} \)}{\( (-\infty;-25) \)}{}{}}
\LANGsk{
\Question{
Vyberte množinu, ktorá je riešením danej nerovnice.
\[
|x-25| \leq 0
\]}
{\( \{25\} \)}{\( \emptyset \)}{\( \mathbb{R} \setminus \{25\} \)}{\( (-\infty;-25) \)}{}{}}
\LANGes{
\Question{
Identifica el conjunto de soluciones de la siguiente inecuación. 
\[
|x-25| \leq 0
\]}
{\( \{25\} \)}{\( \emptyset \)}{\( \mathbb{R} \setminus \{25\} \)}{\( (-\infty;-25) \)}{}{}}
\End

\Data\ID{2000002204}
\Updated{2021-12-01 03:12:11}
\Level{A}
\SubArea{Absolute value equations and inequalities}
\LANGen{
\Question{
Identify the solution set of the following inequality. 
\[
-|2x-100|< 0
\]}
{\( \mathbb{R} \setminus \{50\} \)}{\( \mathbb{R} \)}{\( \mathbb{R} \setminus \{-50\} \)}{\(  \{50\} \)}{}{}}
\LANGpl{
\Question{
Wyznacz zbiór rozwiązań następującej nierówności. 
\[
-|2x-100|< 0
\]}
{\( \mathbb{R} \setminus \{50\} \)}{\( \mathbb{R} \)}{\( \mathbb{R} \setminus \{-50\} \)}{\(  \{50\} \)}{}{}}
\LANGcs{
\Question{
Vyberte množinu, která je řešením dané nerovnice.
\[
-|2x-100|< 0
\]}
{\( \mathbb{R} \setminus \{50\} \)}{\( \mathbb{R} \)}{\( \mathbb{R} \setminus \{-50\} \)}{\(  \{50\} \)}{}{}}
\LANGsk{
\Question{
Vyberte množinu, ktorá je riešením danej nerovnice.
\[
-|2x-100|< 0
\]}
{\( \mathbb{R} \setminus \{50\} \)}{\( \mathbb{R} \)}{\( \mathbb{R} \setminus \{-50\} \)}{\(  \{50\} \)}{}{}}
\LANGes{
\Question{
Identifica el conjunto de soluciones de la siguiente inecuación.  
\[
-|2x-100|< 0
\]}
{\( \mathbb{R} \setminus \{50\} \)}{\( \mathbb{R} \)}{\( \mathbb{R} \setminus \{-50\} \)}{\(  \{50\} \)}{}{}}
\End

\Data\ID{2000002205}
\Updated{2021-12-01 03:12:05}
\Level{A}
\SubArea{Absolute value equations and inequalities}
\LANGen{
\Question{
Identify the solution set of the following inequality.
\[
|7x+1|+1>0
\]}
{\( \mathbb{R} \)}{\( \emptyset \)}{\( \mathbb{R} \setminus \left\{\frac{1}{7}\right\} \)}{\( \mathbb{R} \setminus \left\{-\frac{1}{7}\right\} \)}{}{}}
\LANGpl{
\Question{
Wyznacz zbiór rozwiązań następującej nierówności. 
\[
|7x+1|+1>0
\]}
{\( \mathbb{R} \)}{\( \emptyset \)}{\( \mathbb{R} \setminus \left\{\frac{1}{7}\right\} \)}{\( \mathbb{R} \setminus \left\{-\frac{1}{7}\right\} \)}{}{}}
\LANGcs{
\Question{
Vyberte množinu, která je řešením dané nerovnice.
\[
|7x+1|+1>0
\]}
{\( \mathbb{R} \)}{\( \emptyset \)}{\( \mathbb{R} \setminus \left\{\frac{1}{7}\right\} \)}{\( \mathbb{R} \setminus \left\{-\frac{1}{7}\right\} \)}{}{}}
\LANGsk{
\Question{
Vyberte množinu, ktorá je riešením danej nerovnice.
\[
|7x+1|+1>0
\]}
{\( \mathbb{R} \)}{\( \emptyset \)}{\( \mathbb{R} \setminus \left\{\frac{1}{7}\right\} \)}{\( \mathbb{R} \setminus \left\{-\frac{1}{7}\right\} \)}{}{}}
\LANGes{
\Question{
Identifica el conjunto de soluciones de la siguiente inecuación. 
\[
|7x+1|+1>0
\]}
{\( \mathbb{R} \)}{\( \emptyset \)}{\( \mathbb{R} \setminus \left\{\frac{1}{7}\right\} \)}{\( \mathbb{R} \setminus \left\{-\frac{1}{7}\right\} \)}{}{}}
\End

\Data\ID{2010008503}
\Updated{2022-08-25 10:08:18}
\Level{A}
\SubArea{Absolute value equations and inequalities}
\LANGen{
\Question{
Choose the equation which has only one solution.}
{\( 5-|2x+2|-1=4 \)}{\( 5-|2x+2|+1=5 \)}{\( -|2x+2|+3=4 \)}{\( -|2x+2|+1=0 \)}{}{}}
\LANGpl{
\Question{
Wskaż równanie, które ma tylko jedno rozwiązanie.}
{\( 5-|2x+2|-1=4 \)}{\( 5-|2x+2|+1=5 \)}{\( -|2x+2|+3=4 \)}{\( -|2x+2|+1=0 \)}{}{}}
\LANGcs{
\Question{
Vyberte rovnici, která má právě jedno řešení.}
{\( 5-|2x+2|-1=4 \)}{\( 5-|2x+2|+1=5 \)}{\( -|2x+2|+3=4 \)}{\( -|2x+2|+1=0 \)}{}{}}
\LANGsk{
\Question{
Vyberte rovnicu, ktorá má len jedno riešenie.}
{\( 5-|2x+2|-1=4 \)}{\( 5-|2x+2|+1=5 \)}{\( -|2x+2|+3=4 \)}{\( -|2x+2|+1=0 \)}{}{}}
\LANGes{
\Question{
Elige la ecuación que tiene solo una solución.}
{\( 5-|2x+2|-1=4 \)}{\( 5-|2x+2|+1=5 \)}{\( -|2x+2|+3=4 \)}{\( -|2x+2|+1=0 \)}{}{}}
\End

\Data\ID{2010011701}
\Updated{2022-11-02 09:11:14}
\Level{A}
\SubArea{Absolute value equations and inequalities}
\LANGen{
\Question{
How many natural numbers belong to the solution set of the inequality \( |3x-7| < 5 \)?}
{\(3\)}{\(4\)}{\(2\)}{\(1\)}{}{}}
\LANGpl{
\Question{
Ile liczb należy do zbioru rozwiązań nierówności \( |3x-7| < 5 \)?}
{\(3\)}{\(4\)}{\(2\)}{\(1\)}{}{}}
\LANGcs{
\Question{
Kolik přirozených čísel patří do množiny řešení nerovnice \( |3x-7| < 5 \)?}
{\(3\)}{\(4\)}{\(2\)}{\(1\)}{}{}}
\LANGsk{
\Question{
Koľko prirodzených čísel patrí do množiny riešení nerovnice \( |3x-7| < 5 \)?}
{\(3\)}{\(4\)}{\(2\)}{\(1\)}{}{}}
\LANGes{
\Question{
¿Cuántos números naturales son solución de la inecuación \( |3x-7| < 5 \)?}
{\(3\)}{\(4\)}{\(2\)}{\(1\)}{}{}}
\End

\Data\ID{2010011702}
\Updated{2022-11-02 09:11:14}
\Level{A}
\SubArea{Absolute value equations and inequalities}
\LANGen{
\Question{
Find the solution set of the inequality \( |15-5x| \leq 3 \).}
{\( \left[ \frac{12}5;\frac{18}5\right]\)}{\( \left[ \frac{4}3;\frac{14}3\right]\)}{\( \left[ -\frac{18}5;-\frac{12}5\right]\)}{\( \left[ -\frac{14}3;-\frac{4}3\right]\)}{}{}}
\LANGpl{
\Question{
Wskaż zbiór rozwiązań nierówności \( |15-5x| \leq 3 \).}
{\( \left\langle \frac{12}5;\frac{18}5\right\rangle\)}{\( \left\langle \frac{4}3;\frac{14}3\right\rangle\)}{\( \left\langle -\frac{18}5;-\frac{12}5\right\rangle\)}{\( \left\langle -\frac{14}3;-\frac{4}3\right\rangle\)}{}{}}
\LANGcs{
\Question{
Určete množinu řešení nerovnice \( |15-5x| \leq 3 \).}
{\( \left\langle \frac{12}5;\frac{18}5\right\rangle\)}{\( \left\langle \frac{4}3;\frac{14}3\right\rangle\)}{\( \left\langle -\frac{18}5;-\frac{12}5\right\rangle\)}{\( \left\langle -\frac{14}3;-\frac{4}3\right\rangle\)}{}{}}
\LANGsk{
\Question{
Nájdite množinu riešení nerovnice \( |15-5x| \leq 3 \).}
{\( \left\langle \frac{12}5;\frac{18}5\right\rangle\)}{\( \left\langle \frac{4}3;\frac{14}3\right\rangle\)}{\( \left\langle -\frac{18}5;-\frac{12}5\right\rangle\)}{\( \left\langle -\frac{14}3;-\frac{4}3\right\rangle\)}{}{}}
\LANGes{
\Question{
Halla el conjunto solución de la inecuación \( |15-5x| \leq 3 \).}
{\( \left[ \frac{12}5;\frac{18}5\right]\)}{\( \left[ \frac{4}3;\frac{14}3\right]\)}{\( \left[ -\frac{18}5;-\frac{12}5\right]\)}{\( \left[ -\frac{14}3;-\frac{4}3\right]\)}{}{}}
\End

\Data\ID{2010011704}
\Updated{2022-11-02 09:11:14}
\Level{A}
\SubArea{Absolute value equations and inequalities}
\IMG{2010011701-figure0.pdf}
\LANGen{
\Question{
The solution set of an inequality is graphed on the number line. Determine this inequality.}
{\( |7-x| > 34 \)}{\( |x+7| > 20 \)}{\( |14-x| > 27 \)}{\( |x+14| > 13 \)}{}{}}
\LANGpl{
\Question{
Rozwiązanie nierówności przedstawiono na osi liczbowej. Wskaż tę nierówność.}
{\( |7-x| > 34 \)}{\( |x+7| > 20 \)}{\( |14-x| > 27 \)}{\( |x+14| > 13 \)}{}{}}
\LANGcs{
\Question{
Řešení nerovnice je znázorněno na číselné ose. Určete tuto nerovnici.}
{\( |7-x| > 34 \)}{\( |x+7| > 20 \)}{\( |14-x| > 27 \)}{\( |x+14| > 13 \)}{}{}}
\LANGsk{
\Question{
Riešenie nerovnice je znázornené na číselnej osi. Určte túto nerovnicu.}
{\( |7-x| > 34 \)}{\( |x+7| > 20 \)}{\( |14-x| > 27 \)}{\( |x+14| > 13 \)}{}{}}
\LANGes{
\Question{
El conjunto de soluciones de una inecuación está representado en la recta numérica. Determina esta inecuación.}
{\( |7-x| > 34 \)}{\( |x+7| > 20 \)}{\( |14-x| > 27 \)}{\( |x+14| > 13 \)}{}{}}
\End

\Data\ID{2010011705}
\Updated{2022-11-02 09:11:14}
\Level{A}
\SubArea{Absolute value equations and inequalities}
\IMG{2010011701-figure1.pdf}
\LANGen{
\Question{
The graph shows the set of all real numbers that satisfy the inequality
\( |2x-14| \geq 6\). Determine \( k \).}
{\( k = 4\)}{\( k = 0\)}{\( k = -4\)}{\( k = 6\)}{}{}}
\LANGpl{
\Question{
Wykres przedstawia zbiór wszystkich liczb rzeczywistych spełniających nierówność
\( |2x-14| \geq 6\). Wskaż \( k \).}
{\( k = 4\)}{\( k = 0\)}{\( k = -4\)}{\( k = 6\)}{}{}}
\LANGcs{
\Question{
Graf znázorňuje množinu všech reálných čísel platných pro nerovnici 
\( |2x-14| \geq 6\). Určete \( k \).}
{\( k = 4\)}{\( k = 0\)}{\( k = -4\)}{\( k = 6\)}{}{}}
\LANGsk{
\Question{
Graf znázorňuje množinu všetkých reálnych čísel platných pre nerovnicu 
\( |2x-14| \geq 6\). Určte \( k \).}
{\( k = 4\)}{\( k = 0\)}{\( k = -4\)}{\( k = 6\)}{}{}}
\LANGes{
\Question{
La gráfica muestra el conjunto de todos los números reales que cumplen la inecuación
\( |2x-14| \geq 6\). Determina \( k \).}
{\( k = 4\)}{\( k = 0\)}{\( k = -4\)}{\( k = 6\)}{}{}}
\End

\Data\ID{9000027301}
\Updated{2020-08-23 01:08:35}
\Level{A}
\SubArea{Absolute value equations and inequalities}
\LANGen{
\Question{
Identify the solution set of the following inequality. 
\[ 
 |x| < 3
\]}
{\(\left (-3;3\right )\)}{\(\left (0;3\right )\)}{\(\left [ -3;3\right ] \)}{\(\left [ 0;3\right ] \)}{}{}}
\LANGpl{
\Question{
Wyznacz zbiór rozwiązań podanej nierówności.
\[ 
 |x| < 3
\]}
{\(\left (-3;3\right )\)}{\(\left (0;3\right )\)}{\(\left [ -3;3\right ] \)}{\(\left [ 0;3\right ] \)}{}{}}
\LANGcs{
\Question{
Vyberte množinu, která je řešením nerovnice
\(|x| < 3\).}
{\(\left (-3;3\right )\)}{\(\left (0;3\right )\)}{\(\left \langle -3;3\right \rangle \)}{\(\left \langle 0;3\right \rangle \)}{}{}}
\LANGsk{
\Question{
Vyberte množinu, ktorá je riešením nerovnice.
\[ 
 |x| < 3
\]}
{\(\left (-3;3\right )\)}{\(\left (0;3\right )\)}{\(\left [ -3;3\right ] \)}{\(\left [ 0;3\right ] \)}{}{}}
\LANGes{
\Question{
Identifica el conjunto de soluciones de la siguiente inecuación. 
\[ 
 |x| < 3
\]}
{\(\left (-3;3\right )\)}{\(\left (0;3\right )\)}{\(\left [ -3;3\right ] \)}{\(\left [ 0;3\right ] \)}{}{}}
\End

\Data\ID{9000027302}
\Updated{2020-08-23 01:08:05}
\Level{A}
\SubArea{Absolute value equations and inequalities}
\LANGen{
\Question{
Identify the solution set of the following inequality. 
\[ 
 |x|\geq 5
\]}
{\(\left (-\infty ;-5\right ] \cup \left [ 5;\infty \right )\)}{\(\left (-5;5\right )\)}{\(\left (-\infty ;5\right )\cup \left (5;\infty \right )\)}{\(\left [ -5;\infty \right )\)}{}{}}
\LANGpl{
\Question{
Wyznacz zbiór rozwiązań podanej nierówności.
\[ 
 |x|\geq 5
\]}
{\(\left (-\infty ;-5\right ] \cup \left [ 5;\infty \right )\)}{\(\left (-5;5\right )\)}{\(\left (-\infty ;5\right )\cup \left (5;\infty \right )\)}{\(\left [ -5;\infty \right )\)}{}{}}
\LANGcs{
\Question{
Vyberte množinu, která je řešením nerovnice
\(|x|\geq 5\).}
{\(\left (-\infty ;-5\right \rangle \cup \left \langle 5;\infty \right )\)}{\(\left (-5;5\right )\)}{\(\left (-\infty ;5\right )\cup \left (5;\infty \right )\)}{\(\left \langle -5;\infty \right )\)}{}{}}
\LANGsk{
\Question{
Vyberte množinu, ktorá je riešením danej nerovnice. 
\[ 
 |x|\geq 5
\]}
{\(\left (-\infty ;-5\right ] \cup \left [ 5;\infty \right )\)}{\(\left (-5;5\right )\)}{\(\left (-\infty ;5\right )\cup \left (5;\infty \right )\)}{\(\left [ -5;\infty \right )\)}{}{}}
\LANGes{
\Question{
Identifica el conjunto de soluciones de la siguiente inecuación.
\[ 
 |x|\geq 5
\]}
{\(\left (-\infty ;-5\right ] \cup \left [ 5;\infty \right )\)}{\(\left (-5;5\right )\)}{\(\left (-\infty ;5\right )\cup \left (5;\infty \right )\)}{\(\left [ -5;\infty \right )\)}{}{}}
\End

\Data\ID{9000027303}
\Updated{2020-08-23 01:08:29}
\Level{A}
\SubArea{Absolute value equations and inequalities}
\LANGen{
\Question{
Identify the solution set of the following inequality. 
\[ 
 |x - 4|\leq 1
\]}
{\(\left [ 3;5\right ] \)}{\(\left (-5;-3\right )\)}{\(\left [ -5;-3\right ] \)}{\(\left [ -1;4\right ] \)}{}{}}
\LANGpl{
\Question{
Wyznacz zbiór rozwiązań podanej nierówności.
\[ 
 |x - 4|\leq 1
\]}
{\(\left [ 3;5\right ] \)}{\(\left (-5;-3\right )\)}{\(\left [ -5;-3\right ] \)}{\(\left [ -1;4\right ] \)}{}{}}
\LANGcs{
\Question{
Vyberte množinu, která je řešením nerovnice
\(|x - 4|\leq 1\).}
{\(\left \langle 3;5\right \rangle \)}{\(\left (-5;-3\right )\)}{\(\left \langle -5;-3\right \rangle \)}{\(\left \langle -1;4\right \rangle \)}{}{}}
\LANGsk{
\Question{
Vyberte množinu, ktorá je riešením danej nerovnice. 
\[ 
 |x - 4|\leq 1
\]}
{\(\left [ 3;5\right ] \)}{\(\left (-5;-3\right )\)}{\(\left [ -5;-3\right ] \)}{\(\left [ -1;4\right ] \)}{}{}}
\LANGes{
\Question{
Identifica el conjunto de soluciones de la siguiente inecuación.
\[ 
 |x - 4|\leq 1
\]}
{\(\left [ 3;5\right ] \)}{\(\left (-5;-3\right )\)}{\(\left [ -5;-3\right ] \)}{\(\left [ -1;4\right ] \)}{}{}}
\End

\Data\ID{9000027304}
\Updated{2020-08-23 01:08:53}
\Level{A}
\SubArea{Absolute value equations and inequalities}
\LANGen{
\Question{
Identify the solution set of the following inequality. 
\[ 
 |x - 1| > 10
\]}
{\(\left (-\infty ;-9\right )\cup \left (11;\infty \right )\)}{\(\left (-\infty ;9\right )\cup \left (11;\infty \right )\)}{\(\left [ 11;\infty \right )\)}{\(\left (-\infty ;10\right )\cup \left [ 11;\infty \right )\)}{}{}}
\LANGpl{
\Question{
Wyznacz zbiór rozwiązań podanej nierówności.
\[ 
 |x - 1| > 10
\]}
{\(\left (-\infty ;-9\right )\cup \left (11;\infty \right )\)}{\(\left (-\infty ;9\right )\cup \left (11;\infty \right )\)}{\(\left [ 11;\infty \right )\)}{\(\left (-\infty ;10\right )\cup \left [ 11;\infty \right )\)}{}{}}
\LANGcs{
\Question{
Vyberte množinu, která je řešením dané nerovnice.
\[|x - 1| > 10\]}
{\(\left (-\infty ;-9\right )\cup \left (11;\infty \right )\)}{\(\left (-\infty ;9\right )\cup \left (11;\infty \right )\)}{\(\left \langle 11;\infty \right )\)}{\(\left (-\infty ;10\right )\cup \left \langle 11;\infty \right )\)}{}{}}
\LANGsk{
\Question{
Vyberte množinu, ktorá je riešením danej nerovnice. 
\[ 
 |x - 1| > 10
\]}
{\(\left (-\infty ;-9\right )\cup \left (11;\infty \right )\)}{\(\left (-\infty ;9\right )\cup \left (11;\infty \right )\)}{\(\left [ 11;\infty \right )\)}{\(\left (-\infty ;10\right )\cup \left [ 11;\infty \right )\)}{}{}}
\LANGes{
\Question{
Identifica el conjunto de soluciones de la siguiente inecuación.
\[ 
 |x - 1| > 10
\]}
{\(\left (-\infty ;-9\right )\cup \left (11;\infty \right )\)}{\(\left (-\infty ;9\right )\cup \left (11;\infty \right )\)}{\(\left [ 11;\infty \right )\)}{\(\left (-\infty ;10\right )\cup \left [ 11;\infty \right )\)}{}{}}
\End

\Data\ID{9000027305}
\Updated{2020-08-23 01:08:19}
\Level{A}
\SubArea{Absolute value equations and inequalities}
\LANGen{
\Question{
Identify the solution set of the following inequality. 
\[ 
 |x + 2| < 1
\]}
{\(\left (-3;-1\right )\)}{\(\left (1;3\right )\)}{\(\left [ -1;3\right ] \)}{\(\left [ -2;0\right ] \)}{}{}}
\LANGpl{
\Question{
Wyznacz zbiór rozwiązań podanej nierówności.
\[ 
 |x + 2| < 1
\]}
{\(\left (-3;-1\right )\)}{\(\left (1;3\right )\)}{\(\left [ -1;3\right ] \)}{\(\left [ -2;0\right ] \)}{}{}}
\LANGcs{
\Question{
Vyberte množinu, která je řešením dané nerovnice.
\[|x + 2| < 1\]}
{\(\left (-3;-1\right )\)}{\(\left (1;3\right )\)}{\(\left \langle -1;3\right \rangle \)}{\(\left \langle -2;0\right \rangle \)}{}{}}
\LANGsk{
\Question{
Vyberte množinu, ktorá je riešením danej nerovnice. 
\[ 
 |x + 2| < 1
\]}
{\(\left (-3;-1\right )\)}{\(\left (1;3\right )\)}{\(\left [ -1;3\right ] \)}{\(\left [ -2;0\right ] \)}{}{}}
\LANGes{
\Question{
Identifica el conjunto de soluciones de la siguiente inecuación.
\[ 
 |x + 2| < 1
\]}
{\(\left (-3;-1\right )\)}{\(\left (1;3\right )\)}{\(\left [ -1;3\right ] \)}{\(\left [ -2;0\right ] \)}{}{}}
\End

\Data\ID{9000027306}
\Updated{2020-08-23 01:08:44}
\Level{A}
\SubArea{Absolute value equations and inequalities}
\LANGen{
\Question{
Identify the solution set of the following inequality. 
\[ 
 |x + 3|\geq 6
\]}
{\(\left (-\infty ;-9\right ] \cup \left [ 3;\infty \right )\)}{\(\left (-\infty ;3\right ] \cup \left [ 6;\infty \right )\)}{\(\left (-\infty ;-3\right )\cup \left (9;\infty \right )\)}{\(\left [ -3;6\right ] \)}{}{}}
\LANGpl{
\Question{
Wyznacz zbiór rozwiązań podanej nierówności.
\[ 
 |x + 3|\geq 6
\]}
{\(\left (-\infty ;-9\right ] \cup \left [ 3;\infty \right )\)}{\(\left (-\infty ;3\right ] \cup \left [ 6;\infty \right )\)}{\(\left (-\infty ;-3\right )\cup \left (9;\infty \right )\)}{\(\left [ -3;6\right ] \)}{}{}}
\LANGcs{
\Question{
Vyberte množinu, která je řešením dané nerovnice.
\[|x + 3|\geq 6\]}
{\(\left (-\infty ;-9\right \rangle \cup \left \langle 3;\infty \right )\)}{\(\left (-\infty ;3\right \rangle \cup \left \langle 6;\infty \right )\)}{\(\left (-\infty ;-3\right )\cup \left (9;\infty \right )\)}{\(\left \langle -3;6\right \rangle \)}{}{}}
\LANGsk{
\Question{
Vyberte množinu, ktorá je riešením danej nerovnice. 
\[ 
 |x + 3|\geq 6
\]}
{\(\left (-\infty ;-9\right ] \cup \left [ 3;\infty \right )\)}{\(\left (-\infty ;3\right ] \cup \left [ 6;\infty \right )\)}{\(\left (-\infty ;-3\right )\cup \left (9;\infty \right )\)}{\(\left [ -3;6\right ] \)}{}{}}
\LANGes{
\Question{
Identifica el conjunto de soluciones de la siguiente inecuación.
\[ 
 |x + 3|\geq 6
\]}
{\(\left (-\infty ;-9\right ] \cup \left [ 3;\infty \right )\)}{\(\left (-\infty ;3\right ] \cup \left [ 6;\infty \right )\)}{\(\left (-\infty ;-3\right )\cup \left (9;\infty \right )\)}{\(\left [ -3;6\right ] \)}{}{}}
\End

\Data\ID{9000027307}
\Updated{2020-08-23 12:08:37}
\Level{A}
\SubArea{Absolute value equations and inequalities}
\LANGen{
\Question{
Identify the solution set of the following inequality. 
\[ 
 |2x - 6|\leq 3
\]}
{\(\left [ 1.5;4.5\right ] \)}{\(\left [ 0;6\right ] \)}{\(\left (2;4\right )\)}{\(\left [ -1;5\right ] \)}{}{}}
\LANGpl{
\Question{
Wyznacz zbiór rozwiązań podanej nierówności.
\[ 
 |2x - 6|\leq 3
\]}
{\(\left [ 1.5;4.5\right ] \)}{\(\left [ 0;6\right ] \)}{\(\left (2;4\right )\)}{\(\left [ -1;5\right ] \)}{}{}}
\LANGcs{
\Question{
Vyberte množinu, která je řešením dané nerovnice.
\[|2x - 6|\leq 3\]}
{\(\left \langle 1{,}5;4{,}5\right \rangle \)}{\(\left \langle 0;6\right \rangle \)}{\(\left (2;4\right )\)}{\(\left \langle -1;5\right \rangle \)}{}{}}
\LANGsk{
\Question{
Vyberte množinu, ktorá je riešením danej nerovnice. 
\[ 
 |2x - 6|\leq 3
\]}
{\(\left [ 1.5;4.5\right ] \)}{\(\left [ 0;6\right ] \)}{\(\left (2;4\right )\)}{\(\left [ -1;5\right ] \)}{}{}}
\LANGes{
\Question{
Identifica el conjunto de soluciones de la siguiente inecuación.
\[ 
 |2x - 6|\leq 3
\]}
{\(\left [ 1.5;4.5\right ] \)}{\(\left [ 0;6\right ] \)}{\(\left (2;4\right )\)}{\(\left [ -1;5\right ] \)}{}{}}
\End

\Data\ID{9000027308}
\Updated{2020-08-23 01:08:06}
\Level{A}
\SubArea{Absolute value equations and inequalities}
\LANGen{
\Question{
Identify the solution set of the following inequality. 
\[ 
 |2x - 1| > 5
\]}
{\(\left (-\infty ;-2\right )\cup \left (3;\infty \right )\)}{\(\left (-\infty ;-4.5\right )\cup \left (5.5;\infty \right )\)}{\(\left (1.5;\infty \right )\)}{\(\left (-\infty ;0\right )\cup \left [ 5;\infty \right )\)}{}{}}
\LANGpl{
\Question{
Wyznacz zbiór rozwiązań podanej nierówności.
\[ 
 |2x - 1| > 5
\]}
{\(\left (-\infty ;-2\right )\cup \left (3;\infty \right )\)}{\(\left (-\infty ;-4.5\right )\cup \left (5.5;\infty \right )\)}{\(\left (1.5;\infty \right )\)}{\(\left (-\infty ;0\right )\cup \left [ 5;\infty \right )\)}{}{}}
\LANGcs{
\Question{
Vyberte množinu, která je řešením dané nerovnice.
\[|2x - 1| > 5\]}
{\(\left (-\infty ;-2\right )\cup \left (3;\infty \right )\)}{\(\left (-\infty ;-4{,}5\right )\cup \left (5{,}5;\infty \right )\)}{\(\left (1{,}5;\infty \right )\)}{\(\left (-\infty ;0\right )\cup \left \langle 5;\infty \right )\)}{}{}}
\LANGsk{
\Question{
Vyberte množinu, ktorá je riešením danej nerovnice. 
\[ 
 |2x - 1| > 5
\]}
{\(\left (-\infty ;-2\right )\cup \left (3;\infty \right )\)}{\(\left (-\infty ;-4.5\right )\cup \left (5.5;\infty \right )\)}{\(\left (1.5;\infty \right )\)}{\(\left (-\infty ;0\right )\cup \left [ 5;\infty \right )\)}{}{}}
\LANGes{
\Question{
Identifica el conjunto de soluciones de la siguiente inecuación.  
\[ 
 |2x - 1| > 5
\]}
{\(\left (-\infty ;-2\right )\cup \left (3;\infty \right )\)}{\(\left (-\infty ;-4.5\right )\cup \left (5.5;\infty \right )\)}{\(\left (1.5;\infty \right )\)}{\(\left (-\infty ;0\right )\cup \left [ 5;\infty \right )\)}{}{}}
\End

\Data\ID{9000027309}
\Updated{2020-08-23 12:08:19}
\Level{A}
\SubArea{Absolute value equations and inequalities}
\LANGen{
\Question{
Identify the solution set of the following inequality. 
\[ 
 |3x + 2| < -1
\]}
{\(\emptyset \)}{\(\left (1;3\right )\)}{\(\left [ -1;3\right ] \)}{\(\left [ -2;0\right ] \)}{}{}}
\LANGpl{
\Question{
Wyznacz zbiór rozwiązań podanej nierówności.
\[ 
 |3x + 2| < -1
\]}
{\(\emptyset \)}{\(\left (1;3\right )\)}{\(\left [ -1;3\right ] \)}{\(\left [ -2;0\right ] \)}{}{}}
\LANGcs{
\Question{
Vyberte množinu, která je řešením dané nerovnice.
\[|3x + 2| < -1\]}
{\(\emptyset \)}{\(\left (1;3\right )\)}{\(\left \langle -1;3\right \rangle \)}{\(\left \langle -2;0\right \rangle \)}{}{}}
\LANGsk{
\Question{
Vyberte množinu, ktorá je riešením danej nerovnice. 
\[ 
 |3x + 2| < -1
\]}
{\(\emptyset \)}{\(\left (1;3\right )\)}{\(\left [ -1;3\right ] \)}{\(\left [ -2;0\right ] \)}{}{}}
\LANGes{
\Question{
Identifica el conjunto de soluciones de la siguiente inecuación.
\[ 
 |3x + 2| < -1
\]}
{\(\emptyset \)}{\(\left (1;3\right )\)}{\(\left [ -1;3\right ] \)}{\(\left [ -2;0\right ] \)}{}{}}
\End

\Data\ID{9000027310}
\Updated{2020-08-23 12:08:08}
\Level{A}
\SubArea{Absolute value equations and inequalities}
\LANGen{
\Question{
Identify the solution set of the following inequality. 
\[ 
 |2x + 11| > 0
\]}
{\(\left (-\infty ;-5.5\right )\cup \left (-5.5;\infty \right )\)}{\(\left (-2;11\right )\)}{\(\left (-\infty ;-11\right )\cup \left (11;\infty \right )\)}{\(\emptyset \)}{}{}}
\LANGpl{
\Question{
Wyznacz zbiór rozwiązań podanej nierówności.
\[ 
 |2x + 11| > 0
\]}
{\(\left (-\infty ;-5.5\right )\cup \left (-5.5;\infty \right )\)}{\(\left (-2;11\right )\)}{\(\left (-\infty ;-11\right )\cup \left (11;\infty \right )\)}{\(\emptyset \)}{}{}}
\LANGcs{
\Question{
Vyberte množinu, která je řešením dané nerovnice.
\[|2x + 11| > 0\]}
{\(\left (-\infty ;-5{,}5\right )\cup \left (-5{,}5;\infty \right )\)}{\(\left (-2;11\right )\)}{\(\left (-\infty ;-11\right )\cup \left (11;\infty \right )\)}{\(\emptyset \)}{}{}}
\LANGsk{
\Question{
Vyberte množinu, ktorá je riešením danej nerovnice. 
\[ 
 |2x + 11| > 0
\]}
{\(\left (-\infty ;-5.5\right )\cup \left (-5.5;\infty \right )\)}{\(\left (-2;11\right )\)}{\(\left (-\infty ;-11\right )\cup \left (11;\infty \right )\)}{\(\emptyset \)}{}{}}
\LANGes{
\Question{
Identifica el conjunto de soluciones de la siguiente inecuación.
\[ 
 |2x + 11| > 0
\]}
{\(\left (-\infty ;-5.5\right )\cup \left (-5.5;\infty \right )\)}{\(\left (-2;11\right )\)}{\(\left (-\infty ;-11\right )\cup \left (11;\infty \right )\)}{\(\emptyset \)}{}{}}
\End

\Data\ID{9000078501}
\Updated{2020-12-29 08:12:02}
\Level{A}
\SubArea{Absolute value equations and inequalities}
\LANGen{
\Question{
Write the following set in an interval notation. 
\[ 
 \{x\in \mathbb{R};|x| > 2\}
\]}
{\((-\infty ;-2)\cup (2;\infty )\)}{\([ 2;\infty ] \)}{\((2;\infty )\)}{\((-\infty ;-2] \cup [ 2;\infty )\)}{}{}}
\LANGpl{
\Question{
Zapisz następujący zbiór za pomocą przedziału.
\[ 
 \{x\in \mathbb{R};|x| > 2\}
\]}
{\((-\infty ;-2)\cup (2;\infty )\)}{\([ 2;\infty ] \)}{\((2;\infty )\)}{\((-\infty ;-2] \cup [ 2;\infty )\)}{}{}}
\LANGcs{
\Question{
Vyberte ekvivalentní zápis množiny
\(A\).
\[ 
 A = \{x\in \mathbb{R};|x| > 2\}
\]}
{\((-\infty ;-2)\cup (2;\infty )\)}{\(\langle 2;\infty \rangle \)}{\((2;\infty )\)}{\((-\infty ;-2\rangle \cup \langle 2;\infty )\)}{}{}}
\LANGsk{
\Question{
Vyberte ekvivalentný zápis danej množiny. 
\[ 
 \{x\in \mathbb{R};|x| > 2\}
\]}
{\((-\infty ;-2)\cup (2;\infty )\)}{\([ 2;\infty ] \)}{\((2;\infty )\)}{\((-\infty ;-2] \cup [ 2;\infty )\)}{}{}}
\LANGes{
\Question{
Elige la notación adecuada para el siguiente conjunto: 
\[ 
 \{x\in \mathbb{R};|x| > 2\}
\]}
{\((-\infty ;-2)\cup (2;\infty )\)}{\([ 2;\infty ] \)}{\((2;\infty )\)}{\((-\infty ;-2] \cup [ 2;\infty )\)}{}{}}
\End

\Data\ID{9000078502}
\Updated{2020-12-29 08:12:24}
\Level{A}
\SubArea{Absolute value equations and inequalities}
\LANGen{
\Question{
Write the following set in an interval notation. 
\[ 
 \{x\in \mathbb{R};|x|\leq 4\}
\]}
{\([ - 4;4] \)}{\((-4;4)\)}{\((-\infty ;-4] \)}{\((-\infty ;-4)\)}{}{}}
\LANGpl{
\Question{
Zapisz następujący zbiór za pomocą przedziału.
\[ 
 \{x\in \mathbb{R};|x|\leq 4\}
\]}
{\([ - 4;4] \)}{\((-4;4)\)}{\((-\infty ;-4] \)}{\((-\infty ;-4)\)}{}{}}
\LANGcs{
\Question{
Vyberte ekvivalentní zápis množiny
\(B\).
\[ 
 B = \{x\in \mathbb{R};|x|\leq 4\}
\]}
{\(\langle - 4;4\rangle \)}{\((-4;4)\)}{\((-\infty ;-4\rangle \)}{\((-\infty ;-4)\)}{}{}}
\LANGsk{
\Question{
Vyberte ekvivalentný zápis danej množiny. 
\[ 
 \{x\in \mathbb{R};|x|\leq 4\}
\]}
{\([ - 4;4] \)}{\((-4;4)\)}{\((-\infty ;-4] \)}{\((-\infty ;-4)\)}{}{}}
\LANGes{
\Question{
Elige la notación adecuada para el siguiente conjunto:  
\[ 
 \{x\in \mathbb{R};|x|\leq 4\}
\]}
{\([ - 4;4] \)}{\((-4;4)\)}{\((-\infty ;-4] \)}{\((-\infty ;-4)\)}{}{}}
\End

\Data\ID{9000078503}
\Updated{2020-12-29 08:12:41}
\Level{A}
\SubArea{Absolute value equations and inequalities}
\LANGen{
\Question{
Write the following set in an interval notation. 
\[ 
 \{x\in \mathbb{R};|x - 3|\geq 5\}
\]}
{\((-\infty ;-2] \cup [ 8;\infty )\)}{\((-\infty ;-8] \cup [ 2;\infty )\)}{\([ 2;\infty )\)}{\([ 8;\infty )\)}{}{}}
\LANGpl{
\Question{
Zapisz następujący zbiór za pomocą przedziału.
\[ 
 \{x\in \mathbb{R};|x - 3|\geq 5\}
\]}
{\((-\infty ;-2] \cup [ 8;\infty )\)}{\((-\infty ;-8] \cup [ 2;\infty )\)}{\([ 2;\infty )\)}{\([ 8;\infty )\)}{}{}}
\LANGcs{
\Question{
Vyberte ekvivalentní zápis množiny
\(A\).
\[ 
 A = \{x\in \mathbb{R};|x - 3|\geq 5\}
\]}
{\((-\infty ;-2\rangle \cup \langle 8;\infty )\)}{\((-\infty ;-8\rangle \cup \langle 2;\infty )\)}{\(\langle 2;\infty )\)}{\(\langle 8;\infty )\)}{}{}}
\LANGsk{
\Question{
Vyberte ekvivalentný zápis danej množiny.
\[ 
 \{x\in \mathbb{R};|x - 3|\geq 5\}
\]}
{\((-\infty ;-2] \cup [ 8;\infty )\)}{\((-\infty ;-8] \cup [ 2;\infty )\)}{\([ 2;\infty )\)}{\([ 8;\infty )\)}{}{}}
\LANGes{
\Question{
Elige la notación adecuada para el siguiente conjunto: 
\[ 
 \{x\in \mathbb{R};|x - 3|\geq 5\}
\]}
{\((-\infty ;-2] \cup [ 8;\infty )\)}{\((-\infty ;-8] \cup [ 2;\infty )\)}{\([ 2;\infty )\)}{\([ 8;\infty )\)}{}{}}
\End

\Data\ID{9000078504}
\Updated{2020-12-29 08:12:59}
\Level{A}
\SubArea{Absolute value equations and inequalities}
\LANGen{
\Question{
Write the following set in an interval notation. 
\[ 
 \{x\in \mathbb{R};|x + 10| > 7\}
\]}
{\((-\infty ;-17)\cup (-3;\infty )\)}{\((-\infty ;3)\cup (17;\infty )\)}{\((-3;\infty )\)}{\((17;\infty )\)}{}{}}
\LANGpl{
\Question{
Zapisz następujący zbiór za pomocą przedziału.
\[ 
 \{x\in \mathbb{R};|x + 10| > 7\}
\]}
{\((-\infty ;-17)\cup (-3;\infty )\)}{\((-\infty ;3)\cup (17;\infty )\)}{\((-3;\infty )\)}{\((17;\infty )\)}{}{}}
\LANGcs{
\Question{
Vyberte ekvivalentní zápis množiny
\(B\).
\[ 
 B = \{x\in \mathbb{R};|x + 10| > 7\}
\]}
{\((-\infty ;-17)\cup (-3;\infty )\)}{\((-\infty ;3)\cup (17;\infty )\)}{\((-3;\infty )\)}{\((17;\infty )\)}{}{}}
\LANGsk{
\Question{
Vyberte ekvivalentný zápis danej množiny. 
\[ 
 \{x\in \mathbb{R};|x + 10| > 7\}
\]}
{\((-\infty ;-17)\cup (-3;\infty )\)}{\((-\infty ;3)\cup (17;\infty )\)}{\((-3;\infty )\)}{\((17;\infty )\)}{}{}}
\LANGes{
\Question{
Elige la notación adecuada para el siguiente conjunto: 
\[ 
 \{x\in \mathbb{R};|x + 10| > 7\}
\]}
{\((-\infty ;-17)\cup (-3;\infty )\)}{\((-\infty ;3)\cup (17;\infty )\)}{\((-3;\infty )\)}{\((17;\infty )\)}{}{}}
\End

\Data\ID{9000081401}
\Updated{2020-08-23 12:08:53}
\Level{A}
\SubArea{Absolute value equations and inequalities}
\IMG{000814_01-figure0.pdf}
\LANGen{
\Question{
Find the inequality which describes the set graphed in the picture.}
{\(|x| < 1;\ x\in \mathbb{R}\)}{\(|x - 1| < 0;\ x\in \mathbb{R}\)}{\(|x| > 1;\ x\in \mathbb{R}\)}{\(|x + 1| < 1;\ x\in \mathbb{R}\)}{\(|x - 1| > 0;\ x\in \mathbb{R}\)}{}}
\LANGpl{
\Question{
Z podanych odpowiedzi wybierz nierówność, której zbiór rozwiązań przedstawiono na rysunku poniżej.}
{\(|x| < 1;\ x\in \mathbb{R}\)}{\(|x - 1| < 0;\ x\in \mathbb{R}\)}{\(|x| > 1;\ x\in \mathbb{R}\)}{\(|x + 1| < 1;\ x\in \mathbb{R}\)}{\(|x - 1| > 0;\ x\in \mathbb{R}\)}{}}
\LANGcs{
\Question{
Určete, která z nabídnutých nerovnic má množinu všech řešení
graficky znázorněnou na obrázku.}
{\(|x| < 1;\ x\in \mathbb{R}\)}{\(|x - 1| < 0;\ x\in \mathbb{R}\)}{\(|x| > 1;\ x\in \mathbb{R}\)}{\(|x + 1| < 1;\ x\in \mathbb{R}\)}{\(|x - 1| > 0;\ x\in \mathbb{R}\)}{}}
\LANGsk{
\Question{
Určte, ktorá z ponúknutých nerovníc má množinu všetkých riešení graficky znázornenú na obrázku.}
{\(|x| < 1;\ x\in \mathbb{R}\)}{\(|x - 1| < 0;\ x\in \mathbb{R}\)}{\(|x| > 1;\ x\in \mathbb{R}\)}{\(|x + 1| < 1;\ x\in \mathbb{R}\)}{\(|x - 1| > 0;\ x\in \mathbb{R}\)}{}}
\LANGes{
\Question{
Encuentra la inecuación que describe el conjunto representado en la imagen.}
{\(|x| < 1;\ x\in \mathbb{R}\)}{\(|x - 1| < 0;\ x\in \mathbb{R}\)}{\(|x| > 1;\ x\in \mathbb{R}\)}{\(|x + 1| < 1;\ x\in \mathbb{R}\)}{\(|x - 1| > 0;\ x\in \mathbb{R}\)}{}}
\End

\Data\ID{9000081402}
\Updated{2020-08-23 12:08:22}
\Level{A}
\SubArea{Absolute value equations and inequalities}
\IMG{000814_02-figure0.pdf}
\LANGen{
\Question{
Find the inequality which describes the set graphed in the picture.}
{\(|x - 1| < 2;\ x\in \mathbb{R}\)}{\(|x + 1| < 2;\ x\in \mathbb{R}\)}{\(|x - 1| > 2;\ x\in \mathbb{R}\)}{\(|x + 1| > 2;\ x\in \mathbb{R}\)}{\(|x - 2| > 1;\ x\in \mathbb{R}\)}{}}
\LANGpl{
\Question{
Z podanych odpowiedzi wybierz nierówność, której zbiór rozwiązań przedstawiono na rysunku poniżej.}
{\(|x - 1| < 2;\ x\in \mathbb{R}\)}{\(|x + 1| < 2;\ x\in \mathbb{R}\)}{\(|x - 1| > 2;\ x\in \mathbb{R}\)}{\(|x + 1| > 2;\ x\in \mathbb{R}\)}{\(|x - 2| > 1;\ x\in \mathbb{R}\)}{}}
\LANGcs{
\Question{
Určete, která z nabídnutých nerovnic má množinu všech řešení znázorněnou na obrázku.}
{\(|x - 1| < 2;\ x\in \mathbb{R}\)}{\(|x + 1| < 2;\ x\in \mathbb{R}\)}{\(|x - 1| > 2;\ x\in \mathbb{R}\)}{\(|x + 1| > 2;\ x\in \mathbb{R}\)}{\(|x - 2| > 1;\ x\in \mathbb{R}\)}{}}
\LANGsk{
\Question{
Určte, ktorá z ponúknutých nerovníc má množinu všetkých riešení znázornenú na obrázku.}
{\(|x - 1| < 2;\ x\in \mathbb{R}\)}{\(|x + 1| < 2;\ x\in \mathbb{R}\)}{\(|x - 1| > 2;\ x\in \mathbb{R}\)}{\(|x + 1| > 2;\ x\in \mathbb{R}\)}{\(|x - 2| > 1;\ x\in \mathbb{R}\)}{}}
\LANGes{
\Question{
Encuentra la inecuación que describe el conjunto representado en la imagen.}
{\(|x - 1| < 2;\ x\in \mathbb{R}\)}{\(|x + 1| < 2;\ x\in \mathbb{R}\)}{\(|x - 1| > 2;\ x\in \mathbb{R}\)}{\(|x + 1| > 2;\ x\in \mathbb{R}\)}{\(|x - 2| > 1;\ x\in \mathbb{R}\)}{}}
\End

\Data\ID{9000081403}
\Updated{2020-08-23 12:08:49}
\Level{A}
\SubArea{Absolute value equations and inequalities}
\IMG{000814_03-figure0.pdf}
\LANGen{
\Question{
Find the inequality which describes the set graphed in the picture.}
{\(|x + 1| > 2;\ x\in \mathbb{R}\)}{\(|x - 1| < 2;\ x\in \mathbb{R}\)}{\(|x + 1| < 2;\ x\in \mathbb{R}\)}{\(|x - 1| > 2;\ x\in \mathbb{R}\)}{\(|x - 2| > 1;\ x\in \mathbb{R}\)}{}}
\LANGpl{
\Question{
Z podanych odpowiedzi wybierz nierówność, której zbiór rozwiązań przedstawiono na rysunku poniżej.}
{\(|x + 1| > 2;\ x\in \mathbb{R}\)}{\(|x - 1| < 2;\ x\in \mathbb{R}\)}{\(|x + 1| < 2;\ x\in \mathbb{R}\)}{\(|x - 1| > 2;\ x\in \mathbb{R}\)}{\(|x - 2| > 1;\ x\in \mathbb{R}\)}{}}
\LANGcs{
\Question{
Určete, která z nabídnutých nerovnic má množinu všech řešení znázorněnou na obrázku.}
{\(|x + 1| > 2;\ x\in \mathbb{R}\)}{\(|x - 1| < 2;\ x\in \mathbb{R}\)}{\(|x + 1| < 2;\ x\in \mathbb{R}\)}{\(|x - 1| > 2;\ x\in \mathbb{R}\)}{\(|x - 2| > 1;\ x\in \mathbb{R}\)}{}}
\LANGsk{
\Question{
Určte, ktorá z ponúknutých nerovníc má množinu všetkých riešení znázornenú na obrázku.}
{\(|x + 1| > 2;\ x\in \mathbb{R}\)}{\(|x - 1| < 2;\ x\in \mathbb{R}\)}{\(|x + 1| < 2;\ x\in \mathbb{R}\)}{\(|x - 1| > 2;\ x\in \mathbb{R}\)}{\(|x - 2| > 1;\ x\in \mathbb{R}\)}{}}
\LANGes{
\Question{
Encuentra la inecuación que describe el conjunto representado en la imagen.}
{\(|x + 1| > 2;\ x\in \mathbb{R}\)}{\(|x - 1| < 2;\ x\in \mathbb{R}\)}{\(|x + 1| < 2;\ x\in \mathbb{R}\)}{\(|x - 1| > 2;\ x\in \mathbb{R}\)}{\(|x - 2| > 1;\ x\in \mathbb{R}\)}{}}
\End

\Data\ID{9000081404}
\Updated{2020-08-23 12:08:22}
\Level{A}
\SubArea{Absolute value equations and inequalities}
\IMG{000814_04-figure0.pdf}
\LANGen{
\Question{
Find the inequality which describes the set graphed in the picture.}
{\(|2 + x| > 1;\ x\in \mathbb{R}\)}{\(|2 + x| < 1;\ x\in \mathbb{R}\)}{\(|2 - x| > 1;\ x\in \mathbb{R}\)}{\(|2 - x| < 1;\ x\in \mathbb{R}\)}{\(|1 + x| > 2;\ x\in \mathbb{R}\)}{}}
\LANGpl{
\Question{
Z podanych odpowiedzi wybierz nierówność, której zbiór rozwiązań przedstawiono na rysunku poniżej.}
{\(|2 + x| > 1;\ x\in \mathbb{R}\)}{\(|2 + x| < 1;\ x\in \mathbb{R}\)}{\(|2 - x| > 1;\ x\in \mathbb{R}\)}{\(|2 - x| < 1;\ x\in \mathbb{R}\)}{\(|1 + x| > 2;\ x\in \mathbb{R}\)}{}}
\LANGcs{
\Question{
Určete, která z nabídnutých nerovnic má množinu všech řešení znázorněnou na obrázku.}
{\(|2 + x| > 1;\ x\in \mathbb{R}\)}{\(|2 + x| < 1;\ x\in \mathbb{R}\)}{\(|2 - x| > 1;\ x\in \mathbb{R}\)}{\(|2 - x| < 1;\ x\in \mathbb{R}\)}{\(|1 + x| > 2;\ x\in \mathbb{R}\)}{}}
\LANGsk{
\Question{
Určte, ktorá z ponúknutých nerovníc má množinu všetkých riešení znázornenú na obrázku.}
{\(|2 + x| > 1;\ x\in \mathbb{R}\)}{\(|2 + x| < 1;\ x\in \mathbb{R}\)}{\(|2 - x| > 1;\ x\in \mathbb{R}\)}{\(|2 - x| < 1;\ x\in \mathbb{R}\)}{\(|1 + x| > 2;\ x\in \mathbb{R}\)}{}}
\LANGes{
\Question{
Encuentra la inecuación que describe el conjunto representado en la imagen.}
{\(|2 + x| > 1;\ x\in \mathbb{R}\)}{\(|2 + x| < 1;\ x\in \mathbb{R}\)}{\(|2 - x| > 1;\ x\in \mathbb{R}\)}{\(|2 - x| < 1;\ x\in \mathbb{R}\)}{\(|1 + x| > 2;\ x\in \mathbb{R}\)}{}}
\End

\Data\ID{9000081405}
\Updated{2020-08-23 12:08:47}
\Level{A}
\SubArea{Absolute value equations and inequalities}
\IMG{000814_05-figure0.pdf}
\LANGen{
\Question{
Find the inequality which describes the set graphed in the picture.}
{\(|2 - x| < 1;\ x\in \mathbb{R}\)}{\(|2 + x| > 1;\ x\in \mathbb{R}\)}{\(|2 + x| < 1;\ x\in \mathbb{R}\)}{\(|2 - x| > 1;\ x\in \mathbb{R}\)}{\(|1 - x| < 2;\ x\in \mathbb{R}\)}{}}
\LANGpl{
\Question{
Z podanych odpowiedzi wybierz nierówność, której zbiór rozwiązań przedstawiono na rysunku poniżej.}
{\(|2 - x| < 1;\ x\in \mathbb{R}\)}{\(|2 + x| > 1;\ x\in \mathbb{R}\)}{\(|2 + x| < 1;\ x\in \mathbb{R}\)}{\(|2 - x| > 1;\ x\in \mathbb{R}\)}{\(|1 - x| < 2;\ x\in \mathbb{R}\)}{}}
\LANGcs{
\Question{
Určete, která z nabídnutých nerovnic má množinu všech řešení znázorněnou na obrázku.}
{\(|2 - x| < 1;\ x\in \mathbb{R}\)}{\(|2 + x| > 1;\ x\in \mathbb{R}\)}{\(|2 + x| < 1;\ x\in \mathbb{R}\)}{\(|2 - x| > 1;\ x\in \mathbb{R}\)}{\(|1 - x| < 2;\ x\in \mathbb{R}\)}{}}
\LANGsk{
\Question{
Určte, ktorá z ponúknutých nerovníc má množinu všetkých riešení znázornenú na obrázku.}
{\(|2 - x| < 1;\ x\in \mathbb{R}\)}{\(|2 + x| > 1;\ x\in \mathbb{R}\)}{\(|2 + x| < 1;\ x\in \mathbb{R}\)}{\(|2 - x| > 1;\ x\in \mathbb{R}\)}{\(|1 - x| < 2;\ x\in \mathbb{R}\)}{}}
\LANGes{
\Question{
Encuentra la inecuación que describe el conjunto representado en la imagen.}
{\(|2 - x| < 1;\ x\in \mathbb{R}\)}{\(|2 + x| > 1;\ x\in \mathbb{R}\)}{\(|2 + x| < 1;\ x\in \mathbb{R}\)}{\(|2 - x| > 1;\ x\in \mathbb{R}\)}{\(|1 - x| < 2;\ x\in \mathbb{R}\)}{}}
\End

\Data\ID{1003047101}
\Updated{2022-04-23 05:04:37}
\Level{B}
\SubArea{Absolute value equations and inequalities}
\LANGen{
\Question{
Choose the equation which is, on the interval \( [-2; 3] \), equivalent to
\[ |-x+3|+|x-7|=|x+2|. \]}
{\( (-x+3)-(x-7)=(x+2) \)}{\( -(-x+3)+(x-7)=(x+2) \)}{\( (-x+3)+(x-7)=-(x+2) \)}{\( (-x+3)-(x-7)=-(x+2) \)}{}{}}
\LANGpl{
\Question{
Wybierz równanie, które w przedziale \( \langle-2; 3\rangle \), jest równoważne
\[ |-x+3|+|x-7|=|x+2|. \]}
{\( (-x+3)-(x-7)=(x+2) \)}{\( -(-x+3)+(x-7)=(x+2) \)}{\( (-x+3)+(x-7)=-(x+2) \)}{\( (-x+3)-(x-7)=-(x+2) \)}{}{}}
\LANGcs{
\Question{
Vyberte rovnici, která je na intervalu \( \langle-2; 3\rangle \) ekvivalentní s rovnicí 
\[ |-x+3|+|x-7|=|x+2|. \]}
{\( (-x+3)-(x-7)=(x+2) \)}{\( -(-x+3)+(x-7)=(x+2) \)}{\( (-x+3)+(x-7)=-(x+2) \)}{\( (-x+3)-(x-7)=-(x+2) \)}{}{}}
\LANGsk{
\Question{
Vyberte rovnicu, ktorá je  na intervale \( \langle-2; 3\rangle \), ekvivalentná s rovnicou
\[ |-x+3|+|x-7|=|x+2|. \]}
{\( (-x+3)-(x-7)=(x+2) \)}{\( -(-x+3)+(x-7)=(x+2) \)}{\( (-x+3)+(x-7)=-(x+2) \)}{\( (-x+3)-(x-7)=-(x+2) \)}{}{}}
\LANGes{
\Question{
Elige la ecuación que, en el intervalo \( [-2; 3] \),  es equivalente a 
\[ |-x+3|+|x-7|=|x+2|. \]}
{\( (-x+3)-(x-7)=(x+2) \)}{\( -(-x+3)+(x-7)=(x+2) \)}{\( (-x+3)+(x-7)=-(x+2) \)}{\( (-x+3)-(x-7)=-(x+2) \)}{}{}}
\End

\Data\ID{1003047102}
\Updated{2020-08-23 07:08:15}
\Level{B}
\SubArea{Absolute value equations and inequalities}
\LANGen{
\Question{
Find the sum of the zero points of all the absolute-value expressions of the equation.
\[ |x-3|=|x+9|-|-2+x| \]
(The zero point is the value of \( x \) for which the absolute-value expression equals zero.)}
{\( -4 \)}{\( -8 \)}{\( 14 \)}{\( 10 \)}{}{}}
\LANGpl{
\Question{
Znajdź sumę punktów zerowych wszystkich wyrażeń o bezwzględnej wartości równania.
\[ |x-3|=|x+9|-|-2+x| \]
(Punktem zerowym jest wartość \( x \), dla której wyrażenie wartości bezwzględnej jest równe zero.)}
{\( -4 \)}{\( -8 \)}{\( 14 \)}{\( 10 \)}{}{}}
\LANGcs{
\Question{
Určete součet nulových bodů výrazů v absolutní hodnotě.
\[ |x-3|=|x+9|-|-2+x| \]}
{\( -4 \)}{\( -8 \)}{\( 14 \)}{\( 10 \)}{}{}}
\LANGsk{
\Question{
Určte súčet nulových bodov rovnice.
\[ |x-3|=|x+9|-|-2+x| \]
(Nulový bod  je taká hodnota \( x \), pre ktorú výraz v absolútnej hodnote sa rovná nule.)}
{\( -4 \)}{\( -8 \)}{\( 14 \)}{\( 10 \)}{}{}}
\LANGes{
\Question{
Halla la suma de los puntos cero de todas las expresiones en valores absolutos. 
\[ |x-3|=|x+9|-|-2+x| \]
(El punto cero es el valor de \( x \) para el que la expresión en valor absoluto es igual a cero.)}
{\( -4 \)}{\( -8 \)}{\( 14 \)}{\( 10 \)}{}{}}
\End

\Data\ID{1003047103}
\Updated{2020-08-23 07:08:34}
\Level{B}
\SubArea{Absolute value equations and inequalities}
\LANGen{
\Question{
Find the set of the zero points of all the absolute values from the equation.
\[ |-x+1|=|3x+9|-|2x|\] 
(The zero point is the value of \( x \) for which the absolute value equals zero.)}
{\( \{ -3;0;1 \} \)}{\( \{ -3;-1;0 \} \)}{\( \{ -3;0 \} \)}{\( \{ -3;-1;3 \} \)}{}{}}
\LANGpl{
\Question{
Wyznacz zbiór punktów zerowych wszystkich wartości bezwzględnych z równania.
\[ |-x+1|=|3x+9|-|2x|\] 
Punktem zerowym jest wartość \( x \), dla której wyrażenie wartości bezwzględnej jest równe zero.)}
{\( \{ -3;0;1 \} \)}{\( \{ -3;-1;0 \} \)}{\( \{ -3;0 \} \)}{\( \{ -3;-1;3 \} \)}{}{}}
\LANGcs{
\Question{
Určete množinu nulových bodů výrazů v absolutní hodnotě.
\[ |-x+1|=|3x+9|-|2x|\]}
{\( \{ -3;0;1 \} \)}{\( \{ -3;-1;0 \} \)}{\( \{ -3;0 \} \)}{\( \{ -3;-1;3 \} \)}{}{}}
\LANGsk{
\Question{
Určte množinu nulových bodov rovnice.
\[ |-x+1|=|3x+9|-|2x|\] 
(Nulový bod je taká hodnota \( x \), pre ktorú výraz v absolútnej hodnote sa rovná nule.)}
{\( \{ -3;0;1 \} \)}{\( \{ -3;-1;0 \} \)}{\( \{ -3;0 \} \)}{\( \{ -3;-1;3 \} \)}{}{}}
\LANGes{
\Question{
Halla el conjunto de los puntos cero de todas las expresiones en valor absoluto. 
\[ |-x+1|=|3x+9|-|2x|\] 
(El punto cero es el valor de \( x \) para el que el valor absoluto es igual a cero.)}
{\( \{ -3;0;1 \} \)}{\( \{ -3;-1;0 \} \)}{\( \{ -3;0 \} \)}{\( \{ -3;-1;3 \} \)}{}{}}
\End

\Data\ID{1003047104}
\Updated{2022-04-23 05:04:37}
\Level{B}
\SubArea{Absolute value equations and inequalities}
\LANGen{
\Question{
Decide which of the following intervals ensures that the expressions in all the absolute values of the given equation are positive on this interval. 
\[ |2x+4|+|-x+1|=2x \]}
{\( (-2;1) \)}{\( (-2;0) \)}{\( ( 0;1) \)}{\( (1;\infty) \)}{}{}}
\LANGpl{
\Question{
Zdecyduj, który z podanych przedziałów zapewnia, że wyrażenia we wszystkich wartościach bezwzględnych danego równania są dodatnie w tym przedziale.
\[ |2x+4|+|-x+1|=2x \]}
{\( (-2;1) \)}{\( (-2;0) \)}{\( ( 0;1) \)}{\( (1;\infty) \)}{}{}}
\LANGcs{
\Question{
Určete interval, na kterém jsou všechny výrazy v absolutních hodnotách dané rovnice kladné.
\[ |2x+4|+|-x+1|=2x \]}
{\( (-2;1) \)}{\( (-2;0) \)}{\( ( 0;1) \)}{\( (1;\infty) \)}{}{}}
\LANGsk{
\Question{
Určte interval, na ktorom sú všetky výrazy v absolútnych hodnotách danej rovnice kladné. 
\[ |2x+4|+|-x+1|=2x \]}
{\( (-2;1) \)}{\( (-2;0) \)}{\( ( 0;1) \)}{\( (1;\infty) \)}{}{}}
\LANGes{
\Question{
Define el intervalo en el que todas las expresiones en valor absoluto de la siguiente ecuación son positivas.  
\[ |2x+4|+|-x+1|=2x \]}
{\( (-2;1) \)}{\( (-2;0) \)}{\( ( 0;1) \)}{\( (1;\infty) \)}{}{}}
\End

\Data\ID{1003047105}
\Updated{2020-12-29 08:12:50}
\Level{B}
\SubArea{Absolute value equations and inequalities}
\LANGen{
\Question{
Decide which of the following intervals ensures that the expressions in all the absolute values of the given equation are positive on this interval.
\[ |x+3|+|2x|=|x-1| \]}
{\( (1;\infty) \)}{\( (-3;1) \)}{\( (-3;\infty) \)}{\( (-\infty;-3) \)}{}{}}
\LANGpl{
\Question{
Zdecyduj, który z podanych przedziałów zapewnia, że wyrażenia we wszystkich wartościach bezwzględnych danego równania są dodatnie w tym przedziale.
\[ |x+3|+|2x|=|x-1| \]}
{\( (1;\infty) \)}{\( (-3;1) \)}{\( (-3;\infty) \)}{\( (-\infty;-3) \)}{}{}}
\LANGcs{
\Question{
Určete interval, na kterém jsou všechny výrazy v absolutních hodnotách dané rovnice kladné.
\[ |x+3|+|2x|=|x-1| \]}
{\( (1;\infty) \)}{\( (-3;1) \)}{\( (-3;\infty) \)}{\( (-\infty;-3) \)}{}{}}
\LANGsk{
\Question{
Určte interval, na ktorom sú všetky výrazy v absolútnych hodnotách danej rovnice kladné.
\[ |x+3|+|2x|=|x-1| \]}
{\( (1;\infty) \)}{\( (-3;1) \)}{\( (-3;\infty) \)}{\( (-\infty;-3) \)}{}{}}
\LANGes{
\Question{
Define el intervalo en el que todas las expresiones en  valor absoluto de la siguiente ecuación son positivas. 
\[ |x+3|+|2x|=|x-1| \]}
{\( (1;\infty) \)}{\( (-3;1) \)}{\( (-3;\infty) \)}{\( (-\infty;-3) \)}{}{}}
\End

\Data\ID{1003047106}
\Updated{2020-12-29 08:12:28}
\Level{B}
\SubArea{Absolute value equations and inequalities}
\LANGen{
\Question{
Decide which of the following intervals ensures that the expressions in all the absolute values of the given equation are negative on this interval.
\[ |5-x|-|x+3|=|x+1| \]}
{No such interval exists.}{\( (-\infty;-3) \)}{\( (-\infty;-1) \)}{\( (-1;  5) \)}{}{}}
\LANGpl{
\Question{
Zdecyduj, który z podanych przedziałów zapewnia, że wyrażenia we wszystkich wartościach bezwzględnych danego równania są ujemne w tym przedziale.
\[ |5-x|-|x+3|=|x+1| \]}
{Taki przedział  nie istnieje.}{\( (-\infty;-3) \)}{\( (-\infty;-1) \)}{\( (-1;  5) \)}{}{}}
\LANGcs{
\Question{
Určete interval, na kterém jsou všechny výrazy v absolutních hodnotách dané rovnice záporné.
\[ |5-x|-|x+3|=|x+1| \]}
{Takový interval neexistuje.}{\( (-\infty;-3) \)}{\( (-\infty;-1) \)}{\( (-1;  5) \)}{}{}}
\LANGsk{
\Question{
Určte interval, na ktorom sú všetky výrazy v absolútnych hodnotách záporné.
\[ |5-x|-|x+3|=|x+1| \]}
{Taký interval neexistuje.}{\( (-\infty;-3) \)}{\( (-\infty;-1) \)}{\( (-1;  5) \)}{}{}}
\LANGes{
\Question{
Define el intervalo en el que todas las expresiones en valor absoluto de la siguiente ecuación son negativas. 
\[ |5-x|-|x+3|=|x+1| \]}
{Ese intervalo no existe.}{\( (-\infty;-3) \)}{\( (-\infty;-1) \)}{\( (-1;  5) \)}{}{}}
\End

\Data\ID{1003047107}
\Updated{2020-12-29 08:12:56}
\Level{B}
\SubArea{Absolute value equations and inequalities}
\LANGen{
\Question{
Decide which of the following intervals ensures that the expressions in all the absolute values of the given equation are positive on this interval.
\[ |2x-1|+|2x|=|1-x|-|x+4| \]}
{\( (0.5;1) \)}{\( (-4; 1) \)}{\( (1;\infty) \)}{\( (-4; 0) \)}{}{}}
\LANGpl{
\Question{
Zdecyduj, który z podanych przedziałów zapewnia, że wyrażenia we wszystkich wartościach bezwzględnych danego równania są dodatnie w tym przedziale.
\[ |2x-1|+|2x|=|1-x|-|x+4| \]}
{\( (0{,}5;1) \)}{\( (-4; 1) \)}{\( (1;\infty) \)}{\( (-4; 0) \)}{}{}}
\LANGcs{
\Question{
Určete interval, na kterém jsou všechny výrazy v absolutních hodnotách dané rovnice kladné.
\[ |2x-1|+|2x|=|1-x|-|x+4| \]}
{\( (0{,}5;1) \)}{\( (-4; 1) \)}{\( (1;\infty) \)}{\( (-4; 0) \)}{}{}}
\LANGsk{
\Question{
Určte interval, na ktorom sú všetky výrazy v absolútnych hodnotách kladné.
\[ |2x-1|+|2x|=|1-x|-|x+4| \]}
{\( (0{,}5;1) \)}{\( (-4; 1) \)}{\( (1;\infty) \)}{\( (-4; 0) \)}{}{}}
\LANGes{
\Question{
Define el intervalo en el que todas las expresiones en valor absoluto de la siguiente ecuación son negativas. 
\[ |2x-1|+|2x|=|1-x|-|x+4| \]}
{\( (0.5;1) \)}{\( (-4; 1) \)}{\( (1;\infty) \)}{\( (-4; 0) \)}{}{}}
\End

\Data\ID{1003162901}
\Updated{2020-12-29 08:12:07}
\Level{B}
\SubArea{Absolute value equations and inequalities}
\LANGen{
\Question{
Calculate for which value(s) of the parameter \( t \) the given equation has at least one solution.
\[ |x-5|=t \]}
{\( t\in[0;\infty) \)}{\( t\in[5;\infty) \)}{\( t\in\{0\} \)}{\( t\in(-\infty;5] \)}{}{}}
\LANGpl{
\Question{
Oszacuj dla jakich wartości parametru \( t \)  podane równanie ma przynajmniej jedno rozwiązanie. 
\[ |x-5|=t \]}
{\( t\in\langle0;\infty) \)}{\( t\in\langle5;\infty) \)}{\( t\in\{0\} \)}{\( t\in(-\infty;5\rangle \)}{}{}}
\LANGcs{
\Question{
Určete všechna \( t \), \( t\in\mathbb{R} \), pro která má daná rovnice s neznámou \( x \) alespoň jedno řešení.
\[ |x-5|=t \]}
{\( t\in\langle0;\infty) \)}{\( t\in\langle5;\infty) \)}{\( t\in\{0\} \)}{\( t\in(-\infty;5\rangle \)}{}{}}
\LANGsk{
\Question{
Určte všetky hodnoty parametra \( t \), pre ktoré má daná rovnica aspoň jedno riešenie.
\[ |x-5|=t \]}
{\( t\in\langle0;\infty) \)}{\( t\in\langle5;\infty) \)}{\( t\in\{0\} \)}{\( t\in(-\infty;5\rangle \)}{}{}}
\LANGes{
\Question{
Calcula para qué valor(es) del parámetro \( t \) la siguiente ecuación tiene por lo menos una solución. 
\[ |x-5|=t \]}
{\( t\in[0;\infty) \)}{\( t\in[5;\infty) \)}{\( t\in\{0\} \)}{\( t\in(-\infty;5] \)}{}{}}
\End

\Data\ID{1003162902}
\Updated{2022-04-23 05:04:37}
\Level{B}
\SubArea{Absolute value equations and inequalities}
\LANGen{
\Question{
Calculate for which value(s) of the parameter \( t \) the given equation has no solution.
\[ |x+4|-2=t \]}
{\( t\in(-\infty;-2) \)}{\( t\in(-\infty;-4) \)}{\( t\in(-4;\infty) \)}{\( t\in(-\infty;2) \)}{}{}}
\LANGpl{
\Question{
Oszacuj dla jakich wartości parametru \( t \) podane równanie nie ma rozwiązania.
\[ |x+4|-2=t \]}
{\( t\in(-\infty;-2) \)}{\( t\in(-\infty;-4) \)}{\( t\in(-4;\infty) \)}{\( t\in(-\infty;2) \)}{}{}}
\LANGcs{
\Question{
Určete všechna \( t \), \( t\in\mathbb{R} \), pro která daná rovnice s neznámou \( x \) nemá řešení.
\[ |x+4|-2=t \]}
{\( t\in(-\infty;-2) \)}{\( t\in(-\infty;-4) \)}{\( t\in(-4;\infty) \)}{\( t\in(-\infty;2) \)}{}{}}
\LANGsk{
\Question{
Určte hodnoty parametra \( t \), pre ktoré daná rovnica nemá riešenie.
\[ |x+4|-2=t \]}
{\( t\in(-\infty;-2) \)}{\( t\in(-\infty;-4) \)}{\( t\in(-4;\infty) \)}{\( t\in(-\infty;2) \)}{}{}}
\LANGes{
\Question{
Calcula para qué valor(es) del parámetro \( t \) la siguiente ecuación no tiene solución. 
\[ |x+4|-2=t \]}
{\( t\in(-\infty;-2) \)}{\( t\in(-\infty;-4) \)}{\( t\in(-4;\infty) \)}{\( t\in(-\infty;2) \)}{}{}}
\End

\Data\ID{1003162906}
\Updated{2020-08-23 01:08:31}
\Level{B}
\SubArea{Absolute value equations and inequalities}
\LANGen{
\Question{
Find the number of the roots of the equation \( |x+3|+|x+2|=6 \).}
{\( 2 \)}{\( 1 \)}{\( 4 \)}{\( 3 \)}{}{}}
\LANGpl{
\Question{
Wyznacz liczbę pierwiastków równania \( |x+3|+|x+2|=6 \).}
{\( 2 \)}{\( 1 \)}{\( 4 \)}{\( 3 \)}{}{}}
\LANGcs{
\Question{
Určete počet řešení rovnice \( |x+3|+|x+2|=6 \).}
{\( 2 \)}{\( 1 \)}{\( 4 \)}{\( 3 \)}{}{}}
\LANGsk{
\Question{
Určte počet riešení rovnice \( |x+3|+|x+2|=6 \).}
{\( 2 \)}{\( 1 \)}{\( 4 \)}{\( 3 \)}{}{}}
\LANGes{
\Question{
Halla el número de raíces de la ecuación \( |x+3|+|x+2|=6 \).}
{\( 2 \)}{\( 1 \)}{\( 4 \)}{\( 3 \)}{}{}}
\End

\Data\ID{1003162907}
\Updated{2022-04-23 05:04:37}
\Level{B}
\SubArea{Absolute value equations and inequalities}
\LANGen{
\Question{
Find the sum of all the roots of the equation \( |x+4|=|2x-7| \).}
{\( 12 \)}{\( 11 \)}{\( -11 \)}{\( 23 \)}{}{}}
\LANGpl{
\Question{
Wyznacz sumę wszystkich pierwiastków równania \( |x+4|=|2x-7| \).}
{\( 12 \)}{\( 11 \)}{\( -11 \)}{\( 23 \)}{}{}}
\LANGcs{
\Question{
Určete součet všech kořenů rovnice \( |x+4|=|2x-7| \).}
{\( 12 \)}{\( 11 \)}{\( -11 \)}{\( 23 \)}{}{}}
\LANGsk{
\Question{
Určte súčet všetkých koreňov rovnice \( |x+4|=|2x-7| \).}
{\( 12 \)}{\( 11 \)}{\( -11 \)}{\( 23 \)}{}{}}
\LANGes{
\Question{
Halla la suma de todas las raíces de la ecuación \( |x+4|=|2x-7| \).}
{\( 12 \)}{\( 11 \)}{\( -11 \)}{\( 23 \)}{}{}}
\End

\Data\ID{1003163001}
\Updated{2022-04-23 05:04:37}
\Level{B}
\SubArea{Absolute value equations and inequalities}
\LANGen{
\Question{
Find all \( t \), \( t\in\mathbb{R} \), such that the following equation with the variable \( x \) has exactly two solutions.
\[ |x|+t=-3 \]}
{\( t\in(-\infty;-3) \)}{\( t\in\{-3\} \)}{\( t\in(-3;\infty) \)}{\( t\in\{-3;3\} \)}{}{}}
\LANGpl{
\Question{
Wyznacz wszystkie  \( t \), \( t\in\mathbb{R} \), takie dla których podane równanie ze zmienną \( x \) ma dokładnie dwa rozwiązania.
\[ |x|+t=-3 \]}
{\( t\in(-\infty;-3) \)}{\( t\in\{-3\} \)}{\( t\in(-3;\infty) \)}{\( t\in\{-3;3\} \)}{}{}}
\LANGcs{
\Question{
Určete všechna \( t \), \( t\in\mathbb{R} \), pro která má daná rovnice s neznámou \( x \) právě dvě řešení.
\[ |x|+t=-3 \]}
{\( t\in(-\infty;-3) \)}{\( t\in\{-3\} \)}{\( t\in(-3;\infty) \)}{\( t\in\{-3;3\} \)}{}{}}
\LANGsk{
\Question{
Určte všetky \( t \), \( t\in\mathbb{R} \), pre ktoré má daná rovnica s neznámou \( x \) práve dve riešenia.
\[ |x|+t=-3 \]}
{\( t\in(-\infty;-3) \)}{\( t\in\{-3\} \)}{\( t\in(-3;\infty) \)}{\( t\in\{-3;3\} \)}{}{}}
\LANGes{
\Question{
Encuentra todos los valores de \( t \), \( t\in\mathbb{R} \), para los cuales la siguiente ecuación, siendo \( x \) la incognita, tiene exactamente dos soluciones. 
\[ |x|+t=-3 \]}
{\( t\in(-\infty;-3) \)}{\( t\in\{-3\} \)}{\( t\in(-3;\infty) \)}{\( t\in\{-3;3\} \)}{}{}}
\End

\Data\ID{1003163003}
\Updated{2020-12-29 08:12:56}
\Level{B}
\SubArea{Absolute value equations and inequalities}
\LANGen{
\Question{
How many roots has the given equation for \( x\in(-\infty;-3] \)?
\[ |3-2x|-3|x|=-8 \]}
{\( 1 \)}{\( 2 \)}{\( 0 \)}{\( 3 \)}{}{}}
\LANGpl{
\Question{
Ile pierwiastków posiada podane równanie dla \( x\in(-\infty;-3\rangle \)?
\[ |3-2x|-3|x|=-8 \]}
{\( 1 \)}{\( 2 \)}{\( 0 \)}{\( 3 \)}{}{}}
\LANGcs{
\Question{
Určete počet řešení dané rovnice pro \( x\in(-\infty;-3\rangle \).
\[ |3-2x|-3|x|=-8 \]}
{\( 1 \)}{\( 2 \)}{\( 0 \)}{\( 3 \)}{}{}}
\LANGsk{
\Question{
Určte počet riešení danej rovnice pre \( x\in(-\infty;-3\rangle \). 

\[|3-2x|-3|x|=-8\]}
{\( 1 \)}{\( 2 \)}{\( 0 \)}{\( 3 \)}{}{}}
\LANGes{
\Question{
¿Cuántas raíces tiene la ecuación dada para \( x\in(-\infty;-3] \)?
\[ |3-2x|-3|x|=-8 \]}
{\( 1 \)}{\( 2 \)}{\( 0 \)}{\( 3 \)}{}{}}
\End

\Data\ID{1003163005}
\Updated{2020-12-29 08:12:41}
\Level{B}
\SubArea{Absolute value equations and inequalities}
\LANGen{
\Question{
Find all \( t \), \( t\in\mathbb{R} \), such that the following equation with the variable \( x \) has exactly two solutions.
\[ |x-t|+1=3 \]}
{\( t\in\mathbb{R} \)}{\( t\in(-\infty;2) \)}{\( t\in(2;\infty) \)}{\( t\in(-2;\infty) \)}{}{}}
\LANGpl{
\Question{
Wyznacz wszystkie \( t \), \( t\in\mathbb{R} \), takie, dla których równanie ze zmienną  \( x \) ma dokładnie dwa rozwiązania.
\[ |x-t|+1=3 \]}
{\( t\in\mathbb{R} \)}{\( t\in(-\infty;2) \)}{\( t\in(2;\infty) \)}{\( t\in(-2;\infty) \)}{}{}}
\LANGcs{
\Question{
Určete všechna \( t \), \( t\in\mathbb{R} \), pro která má daná rovnice s neznámou \( x \) právě dvě řešení.
\[ |x-t|+1=3 \]}
{\( t\in\mathbb{R} \)}{\( t\in(-\infty;2) \)}{\( t\in(2;\infty) \)}{\( t\in(-2;\infty) \)}{}{}}
\LANGsk{
\Question{
Určte všetky \( t \), \( t\in\mathbb{R} \), pre ktoré má daná rovnica s neznámou with \( x \) práve dve riešenia.
\[ |x-t|+1=3 \]}
{\( t\in\mathbb{R} \)}{\( t\in(-\infty;2) \)}{\( t\in(2;\infty) \)}{\( t\in(-2;\infty) \)}{}{}}
\LANGes{
\Question{
Halla todos los valores de \( t \), \( t\in\mathbb{R} \) para los que la siguiente ecuación , donde la incógnita es la \( x \) tiene dos soluciones. 
\[ |x-t|+1=3 \]}
{\( t\in\mathbb{R} \)}{\( t\in(-\infty;2) \)}{\( t\in(2;\infty) \)}{\( t\in(-2;\infty) \)}{}{}}
\End

\Data\ID{1003187303}
\Updated{2020-08-23 12:08:04}
\Level{B}
\SubArea{Absolute value equations and inequalities}
\LANGen{
\Question{
Find the solution set of the equation \( |x|=-x \).}
{\( (-\infty;0] \)}{\( (-1;1) \)}{\( \{-4\} \)}{\( (0;+\infty) \)}{}{}}
\LANGpl{
\Question{
Wyznacz zbiór rozwiązań równania \( |x|=-x \).}
{\( (-\infty;0\rangle \)}{\( (-1;1) \)}{\( \{-4\} \)}{\( (0;+\infty) \)}{}{}}
\LANGcs{
\Question{
Určete množinu všech řešení rovnice \( |x|=-x \).}
{\( (-\infty;0\rangle \)}{\( (-1;1) \)}{\( \{-4\} \)}{\( (0;+\infty) \)}{}{}}
\LANGsk{
\Question{
Určte množinu riešení rovnice \( |x|=-x \).}
{\( (-\infty;0\rangle \)}{\( (-1;1) \)}{\( \{-4\} \)}{\( (0;+\infty) \)}{}{}}
\LANGes{
\Question{
Encuentra el conjunto de soluciones de la ecuación \( |x|=-x \).}
{\( (-\infty;0] \)}{\( (-1;1) \)}{\( \{-4\} \)}{\( (0;+\infty) \)}{}{}}
\End

\Data\ID{1003187309}
\Updated{2022-04-23 05:04:37}
\Level{B}
\SubArea{Absolute value equations and inequalities}
\LANGen{
\Question{
Give the number which satisfies the equation \( |6x+4|+4x=0 \).}
{\( x=-2 \)}{\( x=1 \)}{\( x=2 \)}{\( x=-1 \)}{}{}}
\LANGpl{
\Question{
Podaj liczbę, która spełnia równanie \( |6x+4|+4x=0 \).}
{\( x=-2 \)}{\( x=1 \)}{\( x=2 \)}{\( x=-1 \)}{}{}}
\LANGcs{
\Question{
Určete číslo, které je řešením rovnice \( |6x+4|+4x=0 \).}
{\( x=-2 \)}{\( x=1 \)}{\( x=2 \)}{\( x=-1 \)}{}{}}
\LANGsk{
\Question{
Určte číslo, ktoré vyhovuje rovnici \( |6x+4|+4x=0 \).}
{\( x=-2 \)}{\( x=1 \)}{\( x=2 \)}{\( x=-1 \)}{}{}}
\LANGes{
\Question{
Define el número que es solución de la ecuación \( |6x+4|+4x=0 \).}
{\( x=-2 \)}{\( x=1 \)}{\( x=2 \)}{\( x=-1 \)}{}{}}
\End

\Data\ID{1003187402}
\Updated{2020-12-29 08:12:28}
\Level{B}
\SubArea{Absolute value equations and inequalities}
\LANGen{
\Question{
For a given parameter \( a\in\mathbb{R} \) the number \( \pi-\sqrt[3]5-\sqrt2+7 \) represents one of the solutions of the equation \( |2x|=2a^2 \). Which of the following numbers is also the solution of this equation?}
{\( \sqrt2-7+\sqrt[3]5-\pi \)}{\( \sqrt{7-\pi+\sqrt[3]5+\sqrt2} \)}{\( \sqrt[3]5-\pi-\sqrt2-7 \)}{\( \sqrt{\pi-\sqrt[3]5-\sqrt2+7} \)}{}{}}
\LANGpl{
\Question{
Dla podanego parametru \( a\in\mathbb{R} \) liczba \( \pi-\sqrt[3]5-\sqrt2+7 \) stanowi jedno z rozwiązań  równania \( |2x|=2a^2 \). Która z podanych liczb również jest rozwiązaniem tego równania?}
{\( \sqrt2-7+\sqrt[3]5-\pi \)}{\( \sqrt{7-\pi+\sqrt[3]5+\sqrt2} \)}{\( \sqrt[3]5-\pi-\sqrt2-7 \)}{\( \sqrt{\pi-\sqrt[3]5-\sqrt2+7} \)}{}{}}
\LANGcs{
\Question{
Pro daný parametr \( a\in\mathbb{R} \) číslo \( \pi-\sqrt[3]5-\sqrt2+7 \) představuje jedno řešení rovnice \( |2x|=2a^2 \). Které z následujících čísel je také řešením této rovnice?}
{\( \sqrt2-7+\sqrt[3]5-\pi \)}{\( \sqrt{7-\pi+\sqrt[3]5+\sqrt2} \)}{\( \sqrt[3]5-\pi-\sqrt2-7 \)}{\( \sqrt{\pi-\sqrt[3]5-\sqrt2+7} \)}{}{}}
\LANGsk{
\Question{
Pre daný parameter \( a\in\mathbb{R} \) číslo \( \pi-\sqrt[3]5-\sqrt2+7 \) predstavuje jedno riešenie rovnice \( |2x|=2a^2 \). Ktoré z nasledujúcich čísel je tiež riešením tejto rovnice?}
{\( \sqrt2-7+\sqrt[3]5-\pi \)}{\( \sqrt{7-\pi+\sqrt[3]5+\sqrt2} \)}{\( \sqrt[3]5-\pi-\sqrt2-7 \)}{\( \sqrt{\pi-\sqrt[3]5-\sqrt2+7} \)}{}{}}
\LANGes{
\Question{
Dado el parámetro \( a\in\mathbb{R} \) , el número \( \pi-\sqrt[3]5-\sqrt2+7 \) representa una de las soluciones de la ecuación \( |2x|=2a^2 \). ¿Cuál de los siguientes números es también solución de esta ecuación?}
{\( \sqrt2-7+\sqrt[3]5-\pi \)}{\( \sqrt{7-\pi+\sqrt[3]5+\sqrt2} \)}{\( \sqrt[3]5-\pi-\sqrt2-7 \)}{\( \sqrt{\pi-\sqrt[3]5-\sqrt2+7} \)}{}{}}
\End

\Data\ID{1003187411}
\Updated{2022-04-23 05:04:48}
\Level{B}
\SubArea{Absolute value equations and inequalities}
\LANGen{
\Question{
The solution set of the equation  \( 3|51-x|-2|x-81|  =0 \) is:}
{\( \{-9;63\} \)}{\( \{63\} \)}{\( \{9;63\} \)}{\( \{-63;-9\} \)}{}{}}
\LANGpl{
\Question{
Zbiorem rozwiązań równania  \( 3|51-x|-2|x-81|  =0 \) jest:}
{\( \{-9;63\} \)}{\( \{63\} \)}{\( \{9;63\} \)}{\( \{-63;-9\} \)}{}{}}
\LANGcs{
\Question{
Množina řešení rovnice \( 3|51-x|-2|x-81|  =0 \) je:}
{\( \{-9;63\} \)}{\( \{63\} \)}{\( \{9;63\} \)}{\( \{-63;-9\} \)}{}{}}
\LANGsk{
\Question{
Množina riešení rovnice  \( 3|51-x|-2|x-81|  =0 \) je:}
{\( \{-9;63\} \)}{\( \{63\} \)}{\( \{9;63\} \)}{\( \{-63;-9\} \)}{}{}}
\LANGes{
\Question{
El conjunto de soluciones de la ecuación \( 3|51-x|-2|x-81|  =0 \) es:}
{\( \{-9;63\} \)}{\( \{63\} \)}{\( \{9;63\} \)}{\( \{-63;-9\} \)}{}{}}
\End

\Data\ID{2010008501}
\Updated{2022-08-25 10:08:18}
\Level{B}
\SubArea{Absolute value equations and inequalities}
\LANGen{
\Question{
Find all \( t\in\mathbb{R} \), such that the following equation with the variable \( x \) has exactly two solutions.
\[ |x+2|+3=t-1 \]}
{\( t\in(4;\infty) \)}{\( t\in(-2;\infty) \)}{\( t\in\{4\} \)}{\( \mathbb{R}\)}{}{}}
\LANGpl{
\Question{
Znajdź wszystkie \( t\in\mathbb{R} \), dla których podane równanie ze zmienną \( x \) ma więcej niż dwa rozwiązania.
\[ |x+2|+3=t-1 \]}
{\( t\in(4;\infty) \)}{\( t\in(-2;\infty) \)}{\( t\in\{4\} \)}{\( \mathbb{R}\)}{}{}}
\LANGcs{
\Question{
Určete všechna \( t\in\mathbb{R} \), pro která má daná rovnice s neznámou  \( x \) právě dvě řešení.
\[ |x+2|+3=t-1 \]}
{\( t\in(4;\infty) \)}{\( t\in(-2;\infty) \)}{\( t\in\{4\} \)}{\( \mathbb{R}\)}{}{}}
\LANGsk{
\Question{
Určte všetky \( t\in\mathbb{R} \), pre ktoré má daná rovnica s neznámou  \( x \) práve dve riešenia.
\[ |x+2|+3=t-1 \]}
{\( t\in(4;\infty) \)}{\( t\in(-2;\infty) \)}{\( t\in\{4\} \)}{\( \mathbb{R}\)}{}{}}
\LANGes{
\Question{
Encuentra todos los valores de \( t\in\mathbb{R} \),  para los cuales la siguiente ecuación, siendo \( x \) la incognita, tiene exactamente dos soluciones.
\[ |x+2|+3=t-1 \]}
{\( t\in(4;\infty) \)}{\( t\in(-2;\infty) \)}{\( t\in\{4\} \)}{\( \mathbb{R}\)}{}{}}
\End

\Data\ID{2010008504}
\Updated{2022-08-25 10:08:18}
\Level{B}
\SubArea{Absolute value equations and inequalities}
\LANGen{
\Question{
Give the number which satisfies the equation \( |x+2|-1=x+3\).}
{\(-3\)}{\( -2 \)}{\( 0 \)}{\( 1 \)}{}{}}
\LANGpl{
\Question{
Wskaż liczbę spełniającą równanie \( |x+2|-1=x+3\).}
{\(-3\)}{\( -2 \)}{\( 0 \)}{\( 1 \)}{}{}}
\LANGcs{
\Question{
Určete číslo, které je řešením rovnice \( |x+2|-1=x+3\).}
{\(-3\)}{\( -2 \)}{\( 0 \)}{\( 1 \)}{}{}}
\LANGsk{
\Question{
Určte číslo, ktoré je riešením rovnice \( |x+2|-1=x+3\).}
{\(-3\)}{\( -2 \)}{\( 0 \)}{\( 1 \)}{}{}}
\LANGes{
\Question{
Define el número que es solución de la ecuación \( |x+2|-1=x+3\).}
{\(-3\)}{\( -2 \)}{\( 0 \)}{\( 1 \)}{}{}}
\End

\Data\ID{2010008601}
\Updated{2022-11-01 09:11:19}
\Level{B}
\SubArea{Absolute value equations and inequalities}
\LANGen{
\Question{
Choose the equation which is, on the interval \( \left[ \frac12; 2 \right] \), equivalent to
\[ |2x-1|+2|x+5|=1-|2-x|. \]}
{\( (2x-1)+2(x+5)=1-(2-x) \)}{\( (2x-1)+2(x+5)=1+(2-x) \)}{\( (2x-1)-2(x+5)=1+(2-x) \)}{\( -(2x-1)+2(x+5)=1+(2-x) \)}{}{}}
\LANGpl{
\Question{
Wybierz równanie, które w przedziale \( \left\langle \frac12; 2 \right\rangle \), jest odpowiednikiem równania 
\[ |2x-1|+2|x+5|=1-|2-x|. \]}
{\( (2x-1)+2(x+5)=1-(2-x) \)}{\( (2x-1)+2(x+5)=1+(2-x) \)}{\( (2x-1)-2(x+5)=1+(2-x) \)}{\( -(2x-1)+2(x+5)=1+(2-x) \)}{}{}}
\LANGcs{
\Question{
Vyberte rovnici, která je na intervalu \( \left\langle \frac12; 2 \right\rangle \) ekvivalentní s rovnicí
\[ |2x-1|+2|x+5|=1-|2-x|. \]}
{\( (2x-1)+2(x+5)=1-(2-x) \)}{\( (2x-1)+2(x+5)=1+(2-x) \)}{\( (2x-1)-2(x+5)=1+(2-x) \)}{\( -(2x-1)+2(x+5)=1+(2-x) \)}{}{}}
\LANGsk{
\Question{
Vyberte rovnicu, ktorá je na intervale \( \left\langle \frac12; 2 \right\rangle \) ekvivalentná s rovnicou
\[ |2x-1|+2|x+5|=1-|2-x|. \]}
{\( (2x-1)+2(x+5)=1-(2-x) \)}{\( (2x-1)+2(x+5)=1+(2-x) \)}{\( (2x-1)-2(x+5)=1+(2-x) \)}{\( -(2x-1)+2(x+5)=1+(2-x) \)}{}{}}
\LANGes{
\Question{
Elige la ecuación que, en el intervalo \( \left[ \frac12; 2 \right] \), es equivalente a
\[ |2x-1|+2|x+5|=1-|2-x|. \]}
{\( (2x-1)+2(x+5)=1-(2-x) \)}{\( (2x-1)+2(x+5)=1+(2-x) \)}{\( (2x-1)-2(x+5)=1+(2-x) \)}{\( -(2x-1)+2(x+5)=1+(2-x) \)}{}{}}
\End

\Data\ID{2010008602}
\Updated{2022-11-02 09:11:07}
\Level{B}
\SubArea{Absolute value equations and inequalities}
\LANGen{
\Question{
Decide which of the following intervals ensures that the expressions in all the absolute values of the given equation are positive on this interval.
\[|2x-3|+ 2|2-x |= 1-3x\]}
{\( \left(\frac32 ;\infty\right) \)}{\( \left(-\infty;2\right) \)}{\( \left(2 ;\infty\right) \)}{\( \left(-\infty;\frac32\right) \)}{}{}}
\LANGpl{
\Question{
Określ przedział, w którym wszystkie wyrażenia w wartościach bezwzględnych danego równania są dodatnie.
\[|2x-3|+ 2|2-x |= 1-3x\]}
{\( \left(\frac32 ;\infty\right) \)}{\( \left(-\infty;2\right) \)}{\( \left(2 ;\infty\right) \)}{\( \left(-\infty;\frac32\right) \)}{}{}}
\LANGcs{
\Question{
Určete interval, na kterém jsou všechny výrazy v absolutních hodnotách dané rovnice kladné.
\[|2x-3|+ 2|2-x |= 1-3x\]}
{\( \left(\frac32 ;\infty\right) \)}{\( \left(-\infty;2\right) \)}{\( \left(2 ;\infty\right) \)}{\( \left(-\infty;\frac32\right) \)}{}{}}
\LANGsk{
\Question{
Určte interval, na ktorom sú všetky výrazy v absolútnych hodnotách danej rovnice kladné.
\[|2x-3|+ 2|2-x |= 1-3x\]}
{\( \left(\frac32 ;\infty\right) \)}{\( \left(-\infty;2\right) \)}{\( \left(2 ;\infty\right) \)}{\( \left(-\infty;\frac32\right) \)}{}{}}
\LANGes{
\Question{
Decide cuál de los siguientes intervalos asegura que las expresiones de todos los valores absolutos de la ecuación dada son positivas en este intervalo.
\[|2x-3|+ 2|2-x |= 1-3x\]}
{\( \left(\frac32 ;\infty\right) \)}{\( \left(-\infty;2\right) \)}{\( \left(2 ;\infty\right) \)}{\( \left(-\infty;\frac32\right) \)}{}{}}
\End

\Data\ID{2010008603}
\Updated{2022-08-25 10:08:18}
\Level{B}
\SubArea{Absolute value equations and inequalities}
\LANGen{
\Question{
Calculate for which values of the parameter \(t \in\mathbb{R} \) the given equation has no solution.
\[ |2-4x|=1-t\]}
{\( t\in (1 ;\infty) \)}{\( t\in (-\infty;1) \)}{\( t\in \left(\frac12 ;\infty\right) \)}{\( t\in \left(-\infty;\frac12\right) \)}{}{}}
\LANGpl{
\Question{
Oblicz, dla których wartości parametru \(t \in\mathbb{R}\) podane równanie nie ma rozwiązania.
\[ |2-4x|=1-t\]}
{\( (1 ;\infty) \)}{\( (-\infty;1) \)}{\( \left(\frac12 ;\infty\right) \)}{\( \left(-\infty;\frac12\right) \)}{}{}}
\LANGcs{
\Question{
Určete všechna \(t \in\mathbb{R}\), pro která daná rovnice s neznámou \(x\) nemá řešení.
\[ |2-4x|=1-t\]}
{\( t\in (1 ;\infty) \)}{\( t\in (-\infty;1) \)}{\(  t\in\left(\frac12 ;\infty\right) \)}{\( t\in \left(-\infty;\frac12\right) \)}{}{}}
\LANGsk{
\Question{
Určte všetky \(t\in\mathbb{R}\), pre ktoré daná rovnice s neznámou \(x\) nemá riešenie.
\[ |2-4x|=1-t\]}
{\( t\in (1 ;\infty) \)}{\( t\in (-\infty;1) \)}{\( t\in \left(\frac12 ;\infty\right) \)}{\( t\in \left(-\infty;\frac12\right) \)}{}{}}
\LANGes{
\Question{
Calcula para qué valores del parámetro \(t\in\mathbb{R}\)  la siguiente ecuación no tiene solución.
\[ |2-4x|=1-t\]}
{\( t\in (1 ;\infty) \)}{\( t\in (-\infty;1) \)}{\( t\in \left(\frac12 ;\infty\right) \)}{\( t\in \left(-\infty;\frac12\right) \)}{}{}}
\End

\Data\ID{2010008605}
\Updated{2022-08-25 10:08:11}
\Level{B}
\SubArea{Absolute value equations and inequalities}
\LANGen{
\Question{
Find the sum of all the roots of the equation \( |2x+5|=|2-x| \).}
{\( -8\)}{\( -9\)}{\( -7\)}{\( -1 \)}{}{}}
\LANGpl{
\Question{
Znajdź sumę wszystkich pierwiastków równania  \( |2x+5|=|2-x| \).}
{\( -8\)}{\( -9\)}{\( -7\)}{\( -1 \)}{}{}}
\LANGcs{
\Question{
Určete součet všech kořenů této rovnice \( |2x+5|=|2-x| \).}
{\( -8\)}{\( -9\)}{\( -7\)}{\( -1 \)}{}{}}
\LANGsk{
\Question{
Určte súčet všetkých koreňov tejto rovnice \( |2x+5|=|2-x| \).}
{\( -8\)}{\( -9\)}{\( -7\)}{\( -1 \)}{}{}}
\LANGes{
\Question{
Halla la suma de todas las raíces de la ecuación \( |2x+5|=|2-x| \).}
{\( -8\)}{\( -9\)}{\( -7\)}{\( -1 \)}{}{}}
\End

\Data\ID{2010011703}
\Updated{2022-11-02 09:11:14}
\Level{B}
\SubArea{Absolute value equations and inequalities}
\LANGen{
\Question{
The solution set of the equation  \( 2|45-x|-|2x-72|  =0 \) is:}
{\( \left\{ \frac{81}2 \right\}\)}{\( \emptyset\)}{\( \left\{ -\frac{81}2 \right\}\)}{\( \left\{ -\frac{81}2 ;\frac{81}2 \right\}\)}{}{}}
\LANGpl{
\Question{
Wskaż zbiór rozwiązań równania  \( 2|45-x|-|2x-72|  =0 \).}
{\( \left\{ \frac{81}2 \right\}\)}{\( \emptyset\)}{\( \left\{ -\frac{81}2 \right\}\)}{\( \left\{ -\frac{81}2 ;\frac{81}2 \right\}\)}{}{}}
\LANGcs{
\Question{
Množina řešení rovnice   \( 2|45-x|-|2x-72|  =0 \) je:}
{\( \left\{ \frac{81}2 \right\}\)}{\( \emptyset\)}{\( \left\{ -\frac{81}2 \right\}\)}{\( \left\{ -\frac{81}2 ;\frac{81}2 \right\}\)}{}{}}
\LANGsk{
\Question{
Množina riešení rovnice \( 2|45-x|-|2x-72|  =0 \) je:}
{\( \left\{ \frac{81}2 \right\}\)}{\( \emptyset\)}{\( \left\{ -\frac{81}2 \right\}\)}{\( \left\{ -\frac{81}2 ;\frac{81}2 \right\}\)}{}{}}
\LANGes{
\Question{
El conjunto solución de la ecuación \( 2|45-x|-|2x-72|  =0 \) es:}
{\( \left\{ \frac{81}2 \right\}\)}{\( \emptyset\)}{\( \left\{ -\frac{81}2 \right\}\)}{\( \left\{ -\frac{81}2 ;\frac{81}2 \right\}\)}{}{}}
\End

\Data\ID{9000026401}
\Updated{2020-08-23 07:08:46}
\Level{B}
\SubArea{Absolute value equations and inequalities}
\LANGen{
\Question{
To solve the equation with absolute value using interval method we have to divide
the domain of the equation by the zero point of the subexpression in the absolute
value. Find this point. 
\[ 
 2x - 1 = 1 + |x|
\]}
{\(0\)}{\(\frac{1}
{2}\)}{\(-\frac{1} 
{2}\)}{\(- 1\)}{}{}}
\LANGpl{
\Question{
Aby rozwiązać równanie z wartością bezwzględną za pomocą metody interwałowej, musimy podzielić dziedzinę tego równania przez punkt zerowy podwyrażenia w wartości bezwzględnej. Znajdź ten punkt.
\[ 
 2x - 1 = 1 + |x|
\]}
{\(0\)}{\(\frac{1}
{2}\)}{\(-\frac{1} 
{2}\)}{\(- 1\)}{}{}}
\LANGcs{
\Question{
Určete nulový bod výrazu v absolutní hodnotě. 
\[ 
 2x - 1 = 1 + |x|
\]}
{\(0\)}{\(\frac{1}
{2}\)}{\(-\frac{1} 
{2}\)}{\(- 1\)}{}{}}
\LANGsk{
\Question{
Určte nulový bod výrazu v absolútnej hodnote.
\[ 
 2x - 1 = 1 + |x|
\]}
{\(0\)}{\(\frac{1}
{2}\)}{\(-\frac{1} 
{2}\)}{\(- 1\)}{}{}}
\LANGes{
\Question{
Encuentra el punto cero de la expresión en el valor absoluto.
\[ 
 2x - 1 = 1 + |x|
\]}
{\(0\)}{\(\frac{1}
{2}\)}{\(-\frac{1} 
{2}\)}{\(- 1\)}{}{}}
\End

\Data\ID{9000026402}
\Updated{2020-08-23 07:08:11}
\Level{B}
\SubArea{Absolute value equations and inequalities}
\LANGen{
\Question{
To solve the equation with absolute value using interval method we have to divide
the domain of the equation by the zero point of the subexpression in the absolute
value. Find this point. 
\[ 
 1 -|x - 2| = x + 2
\]}
{\(2\)}{\(1\)}{\(- 2\)}{\(0\)}{}{}}
\LANGpl{
\Question{
Aby rozwiązać równanie z wartością bezwzględną za pomocą metody interwałowej, musimy podzielić dziedzinę tego równania przez punkt zerowy podwyrażenia w wartości bezwzględnej. Znajdź ten punkt.
\[ 
 1 -|x - 2| = x + 2
\]}
{\(2\)}{\(1\)}{\(- 2\)}{\(0\)}{}{}}
\LANGcs{
\Question{
Určete nulový bod výrazu v absolutní hodnotě. 
\[ 
 1 -|x - 2| = x + 2
\]}
{\(2\)}{\(1\)}{\(- 2\)}{\(0\)}{}{}}
\LANGsk{
\Question{
Určte nulový bod výrazu v absolútnej hodnote. 
\[ 
 1 -|x - 2| = x + 2
\]}
{\(2\)}{\(1\)}{\(- 2\)}{\(0\)}{}{}}
\LANGes{
\Question{
Encuentra el punto cero de la expresión en el valor absoluto.
\[ 
 1 -|x - 2| = x + 2
\]}
{\(2\)}{\(1\)}{\(- 2\)}{\(0\)}{}{}}
\End

\Data\ID{9000026403}
\Updated{2020-08-23 07:08:47}
\Level{B}
\SubArea{Absolute value equations and inequalities}
\LANGen{
\Question{
To solve the equation with absolute value using interval method we have to divide
the domain of the equation by the zero points of the subexpressions in the absolute
value. Find all these points. 
\[ 
 |x + 1| + |2x - 1| = 3
\]}
{\(-1,\ \frac{1} 
{2}\)}{\(- 3\)}{\(1,\ -\frac{1}
{2}\)}{\(0\)}{}{}}
\LANGpl{
\Question{
Aby rozwiązać równanie z wartością bezwzględną za pomocą metody interwałowej, musimy podzielić dziedzinę tego równania przez punkt zerowy podwyrażenia w wartości bezwzględnej. Znajdź ten punkt. 
\[ 
 |x + 1| + |2x - 1| = 3
\]}
{\(-1,\ \frac{1} 
{2}\)}{\(- 3\)}{\(1,\ -\frac{1}
{2}\)}{\(0\)}{}{}}
\LANGcs{
\Question{
Určete nulové body výrazů v absolutní hodnotě. 
\[ 
 |x + 1| + |2x - 1| = 3
\]}
{\(-1,\ \frac{1} 
{2}\)}{\(- 3\)}{\(1,\ -\frac{1}
{2}\)}{\(0\)}{}{}}
\LANGsk{
\Question{
Určte nulové body výrazov v absolútnej hodnote. 
\[ 
 |x + 1| + |2x - 1| = 3
\]}
{\(-1,\ \frac{1} 
{2}\)}{\(- 3\)}{\(1,\ -\frac{1}
{2}\)}{\(0\)}{}{}}
\LANGes{
\Question{
Encuentra todos los puntos cero de las expresiones en el valor absoluto.
\[ 
 |x + 1| + |2x - 1| = 3
\]}
{\(-1,\ \frac{1} 
{2}\)}{\(- 3\)}{\(1,\ -\frac{1}
{2}\)}{\(0\)}{}{}}
\End

\Data\ID{9000026404}
\Updated{2020-08-23 07:08:43}
\Level{B}
\SubArea{Absolute value equations and inequalities}
\LANGen{
\Question{
To solve the equation with absolute value using interval method we have to divide
the domain of the equation by the zero points of the subexpressions in the absolute
value. Find all these points. 
\[ 
 2|x - 2| + |2 - x| = 1 + |x|
\]}
{\(2,\ 0\)}{\(-2,\ 2,\ 0\)}{\(-1,\ 2\)}{\(-1,\ 2,\ 0\)}{}{}}
\LANGpl{
\Question{
Aby rozwiązać równanie z wartością bezwzględną za pomocą metody interwałowej, musimy podzielić dziedzinę tego równania przez punkty zerowy podwyrażeń w wartości bezwzględnej. Znajdź wszystkie te punkty.
\[ 
 2|x - 2| + |2 - x| = 1 + |x|
\]}
{\(2,\ 0\)}{\(-2,\ 2,\ 0\)}{\(-1,\ 2\)}{\(-1,\ 2,\ 0\)}{}{}}
\LANGcs{
\Question{
Určete nulové body výrazů v absolutní hodnotě. 
\[ 
 2|x - 2| + |2 - x| = 1 + |x|
\]}
{\(2,\ 0\)}{\(-2,\ 2,\ 0\)}{\(-1,\ 2\)}{\(-1,\ 2,\ 0\)}{}{}}
\LANGsk{
\Question{
Určte nulové body výrazov v absolútnej hodnote. 
\[ 
 2|x - 2| + |2 - x| = 1 + |x|
\]}
{\(2,\ 0\)}{\(-2,\ 2,\ 0\)}{\(-1,\ 2\)}{\(-1,\ 2,\ 0\)}{}{}}
\LANGes{
\Question{
Encuentra todos los puntos cero de las expresiones en el valor absoluto.
\[ 
 2|x - 2| + |2 - x| = 1 + |x|
\]}
{\(2,\ 0\)}{\(-2,\ 2,\ 0\)}{\(-1,\ 2\)}{\(-1,\ 2,\ 0\)}{}{}}
\End

\Data\ID{9000026405}
\Updated{2020-12-29 08:12:52}
\Level{B}
\SubArea{Absolute value equations and inequalities}
\LANGen{
\Question{
Assuming \(x\in (-\infty ;1)\),
write the following equation in the form which does not contain an absolute value.
\[ 
 3 = |x - 1|
\]}
{\(3 = -x + 1\)}{\(3 = x - 1\)}{\(3 = -x - 1\)}{\(3 = x + 1\)}{}{}}
\LANGpl{
\Question{
Zakładając, że \(x\in (-\infty ;1)\),
zapisz następujące równanie tak, by nie zawierało wartości bezwzględnej.
\[ 
 3 = |x - 1|
\]}
{\(3 = -x + 1\)}{\(3 = x - 1\)}{\(3 = -x - 1\)}{\(3 = x + 1\)}{}{}}
\LANGcs{
\Question{
Rovnici 
\[ 
 3 = |x - 1|
\]
lze v intervalu \((-\infty ;1)\) 
přepsat do tvaru:}
{\(3 = -x + 1\)}{\(3 = x - 1\)}{\(3 = -x - 1\)}{\(3 = x + 1\)}{}{}}
\LANGsk{
\Question{
V intervale \(x\in (-\infty ;1)\) možno rovnicu
\[ 
 3 = |x - 1|
\] prepísať v tvare:}
{\(3 = -x + 1\)}{\(3 = x - 1\)}{\(3 = -x - 1\)}{\(3 = x + 1\)}{}{}}
\LANGes{
\Question{
Sabiendo que \(x\in (-\infty ;1)\),
escribe la siguiente ecuación para que no contenga el valor absoluto. 
\[ 
 3 = |x - 1|
\]}
{\(3 = -x + 1\)}{\(3 = x - 1\)}{\(3 = -x - 1\)}{\(3 = x + 1\)}{}{}}
\End

\Data\ID{9000026406}
\Updated{2020-12-29 08:12:24}
\Level{B}
\SubArea{Absolute value equations and inequalities}
\LANGen{
\Question{
Assuming \(x\in (-3;2)\),
write the following equation in the form which does not contain an absolute value.
\[ 
 |x + 3| = |x - 2|
\]}
{\(x + 3 = -x + 2\)}{\(x + 3 = x - 2\)}{\(- x - 3 = -x + 2\)}{\(- x - 3 = x + 2\)}{}{}}
\LANGpl{
\Question{
Zakładając, że \(x\in (-3;2)\),
zapisz następujące równanie tak, by nie zawierało wartości bezwzględnej. 
\[ 
 |x + 3| = |x - 2|
\]}
{\(x + 3 = -x + 2\)}{\(x + 3 = x - 2\)}{\(- x - 3 = -x + 2\)}{\(- x - 3 = x + 2\)}{}{}}
\LANGcs{
\Question{
Rovnici 
\[ 
 |x + 3| = |x - 2|
\]
lze v intervalu \((-3;2)\) 
přepsat do tvaru:}
{\(x + 3 = -x + 2\)}{\(x + 3 = x - 2\)}{\(- x - 3 = -x + 2\)}{\(- x - 3 = x + 2\)}{}{}}
\LANGsk{
\Question{
V intervale \(x\in (-3;2)\)
možno rovnicu 
\[ 
 |x + 3| = |x - 2|
\] prepísať v tvare:}
{\(x + 3 = -x + 2\)}{\(x + 3 = x - 2\)}{\(- x - 3 = -x + 2\)}{\(- x - 3 = x + 2\)}{}{}}
\LANGes{
\Question{
Sabiendo que \(x\in (-3;2)\),
escribe la forma de la siguiente ecuación sin que contenga un valor absoluto. 
\[ 
 |x + 3| = |x - 2|
\]}
{\(x + 3 = -x + 2\)}{\(x + 3 = x - 2\)}{\(- x - 3 = -x + 2\)}{\(- x - 3 = x + 2\)}{}{}}
\End

\Data\ID{9000026407}
\Updated{2020-12-29 08:12:45}
\Level{B}
\SubArea{Absolute value equations and inequalities}
\LANGen{
\Question{
Assuming \(x\in \left [ \frac{1}
{2};\infty \right )\),
write the following equation in the form which does not contain an absolute value.
\[ 
 |x + 1| + |2x - 1| = 3
\]}
{\(x + 1 + 2x - 1 = 3\)}{\(x + 1 - 2x - 1 = 3\)}{\(x + 1 - 2x + 1 = 3\)}{\(- x - 1 + 2x - 1 = 3\)}{}{}}
\LANGpl{
\Question{
Zakładając, że \(x\in \left [ \frac{1}
{2};\infty \right )\),
zapisz następujące równanie tak, by nie zawierało wartości bezwzględnej. 
\[ 
 |x + 1| + |2x - 1| = 3
\]}
{\(x + 1 + 2x - 1 = 3\)}{\(x + 1 - 2x - 1 = 3\)}{\(x + 1 - 2x + 1 = 3\)}{\(- x - 1 + 2x - 1 = 3\)}{}{}}
\LANGcs{
\Question{
Rovnici 
\[ 
 |x + 1| + |2x - 1| = 3
\]
lze v intervalu \(\left \langle \frac{1}
{2};\infty \right )\) 
přepsat do tvaru:}
{\(x + 1 + 2x - 1 = 3\)}{\(x + 1 - 2x - 1 = 3\)}{\(x + 1 - 2x + 1 = 3\)}{\(- x - 1 + 2x - 1 = 3\)}{}{}}
\LANGsk{
\Question{
V intervale \(x\in \left [ \frac{1}
{2};\infty \right )\),
možno rovnicu 
\[ 
 |x + 1| + |2x - 1| = 3
\] prepísať v tvare:}
{\(x + 1 + 2x - 1 = 3\)}{\(x + 1 - 2x - 1 = 3\)}{\(x + 1 - 2x + 1 = 3\)}{\(- x - 1 + 2x - 1 = 3\)}{}{}}
\LANGes{
\Question{
Sabiendo que \(x\in \left [ \frac{1}
{2};\infty \right )\),
escribe la siguiente ecuación para que no contenga el valor absoluto. 
\[ 
 |x + 1| + |2x - 1| = 3
\]}
{\(x + 1 + 2x - 1 = 3\)}{\(x + 1 - 2x - 1 = 3\)}{\(x + 1 - 2x + 1 = 3\)}{\(- x - 1 + 2x - 1 = 3\)}{}{}}
\End

\Data\ID{9000026409}
\Updated{2020-12-29 08:12:57}
\Level{B}
\SubArea{Absolute value equations and inequalities}
\LANGen{
\Question{
Consider the following equation. 
\[ 
 |2x - 4| = 5x - 7
\]
Solving the equation on the intervals where it is possible to evaluate the absolute
value we get equations on partial subintervals as follows.
\[\begin{aligned}
 \text{for }x &\in (-\infty ;2)\colon &\text{for }x &\in [ 2;\infty )\colon & & & &
 \\ - 2x + 4 & = 5x - 7 &2x - 4 & = 5x - 7 & & & &
 \\ - 7x & = -11 & - 3x & = -3 & & & &
 \\x & = \frac{11} 
 {7} &x & = 1 & & & &
\end{aligned}\]
Find the solution set of the original equation.}
{\(\left \{\frac{11}
 {7} \right \}\)}{\(\left \{\frac{11}
 {7} ;1\right \}\)}{\(\left \{1\right \}\)}{\(\emptyset \)}{}{}}
\LANGpl{
\Question{
Rozważ następujące równanie.
\[ 
 |2x - 4| = 5x - 7
\]
Rozwiązując równania w przedziałach, w których możliwe jest oszacowanie 
wartości bezwzględnej,  otrzymujemy równania na podprzedziałach cząstkowych w następujący sposób.
\[\begin{aligned}
 \text{pro }x &\in (-\infty ;2)\colon &\text{pro }x &\in [ 2;\infty )\colon & & & &
 \\ - 2x + 4 & = 5x - 7 &2x - 4 & = 5x - 7 & & & &
 \\ - 7x & = -11 & - 3x & = -3 & & & &
 \\x & = \frac{11} 
 {7} &x & = 1 & & & &
\end{aligned}\]
Wyznacz zbiór rozwiązań pierwotnego równania.}
{\(\left \{\frac{11}
 {7} \right \}\)}{\(\left \{\frac{11}
 {7} ;1\right \}\)}{\(\left \{1\right \}\)}{\(\emptyset \)}{}{}}
\LANGcs{
\Question{
Nulový bod výrazu v absolutní hodnotě v rovnici 
\[ 
 |2x - 4| = 5x - 7
\]
je \(2\).
Přepsáním pro jednotlivé intervaly dostaneme rovnici a dílčí řešení:
\[\begin{aligned}
 \text{pro }x &\in (-\infty ;2)\colon &\text{pro }x &\in \langle 2;\infty )\colon & & & &
 \\ - 2x + 4 & = 5x - 7 &2x - 4 & = 5x - 7 & & & &
 \\ - 7x & = -11 & - 3x & = -3 & & & &
 \\x & = \frac{11} 
 {7} &x & = 1 & & & &
\end{aligned}\]
Označte správnou množinu kořenů původní rovnice:}
{\(\left \{\frac{11}
 {7} \right \}\)}{\(\left \{\frac{11}
 {7} ;1\right \}\)}{\(\left \{1\right \}\)}{\(\emptyset \)}{}{}}
\LANGsk{
\Question{
Nulový bod výrazu v absolútnej hodnote v rovnici
\[ 
 |2x - 4| = 5x - 7
\] je \(2\).
Prepísaním pre jednotlivé intervaly dostaneme rovnicu a čiastočné riešenia:
\[\begin{aligned}
 \text{pre }x &\in (-\infty ;2)\colon &\text{pre }x &\in [ 2;\infty )\colon & & & &
 \\ - 2x + 4 & = 5x - 7 &2x - 4 & = 5x - 7 & & & &
 \\ - 7x & = -11 & - 3x & = -3 & & & &
 \\x & = \frac{11} 
 {7} &x & = 1 & & & &
\end{aligned}\]
Označte správnu množinu koreňov pôvodnej rovnice.}
{\(\left \{\frac{11}
 {7} \right \}\)}{\(\left \{\frac{11}
 {7} ;1\right \}\)}{\(\left \{1\right \}\)}{\(\emptyset \)}{}{}}
\LANGes{
\Question{
Consideremos la siguiente ecuación. 
\[ 
 |2x - 4| = 5x - 7
\]
Resolviendo la ecuación en los diferentes intervalos donde es posible evaluar el valor absoluto, obtenemos dos ecuaciones en los intervalos parciales siguientes: 
\[\begin{aligned}
 \text{para }x &\in (-\infty ;2)\colon &\text{para }x &\in [ 2;\infty )\colon & & & &
 \\ - 2x + 4 & = 5x - 7 &2x - 4 & = 5x - 7 & & & &
 \\ - 7x & = -11 & - 3x & = -3 & & & &
 \\x & = \frac{11} 
 {7} &x & = 1 & & & &
\end{aligned}\]
Encuentra el conjunto de soluciones de la ecuación original.}
{\(\left \{\frac{11}
 {7} \right \}\)}{\(\left \{\frac{11}
 {7} ;1\right \}\)}{\(\left \{1\right \}\)}{\(\emptyset \)}{}{}}
\End

\Data\ID{1003162903}
\Updated{2022-04-23 05:04:37}
\Level{C}
\SubArea{Absolute value equations and inequalities}
\LANGen{
\Question{
Find the sum of all the roots of the given equation.
\[ \bigl| |x-4|-2\bigr|=0 \]}
{\( 8 \)}{\( 12 \)}{\( 4 \)}{\( 6 \)}{}{}}
\LANGpl{
\Question{
Wyznacz sumę wszystkich pierwiastków podanego równania. 
\[ \bigl| |x-4|-2\bigr|=0 \]}
{\( 8 \)}{\( 12 \)}{\( 4 \)}{\( 6 \)}{}{}}
\LANGcs{
\Question{
Určete součet všech kořenů dané rovnice.
\[ \bigl| |x-4|-2\bigr|=0 \]}
{\( 8 \)}{\( 12 \)}{\( 4 \)}{\( 6 \)}{}{}}
\LANGsk{
\Question{
Určte súčet všetkých koreňov danej rovnice.
\[ \bigl| |x-4|-2\bigr|=0 \]}
{\( 8 \)}{\( 12 \)}{\( 4 \)}{\( 6 \)}{}{}}
\LANGes{
\Question{
Halla la suma de todas las raíces de la ecuación 
\[ \bigl| |x-4|-2\bigr|=0 \].}
{\( 8 \)}{\( 12 \)}{\( 4 \)}{\( 6 \)}{}{}}
\End

\Data\ID{1003162904}
\Updated{2020-12-29 08:12:43}
\Level{C}
\SubArea{Absolute value equations and inequalities}
\LANGen{
\Question{
Find the sum of all the roots of the equation \( \bigl| |1-x|-2 \bigr|=1 \) for \( x\in(1;\infty) \).}
{\( 6 \)}{\( 4 \)}{\( 1 \)}{\( 2 \)}{}{}}
\LANGpl{
\Question{
Wyznacz sumę wszystkich pierwiastków równania \( \bigl| |1-x|-2 \bigr|=1 \) pro \( x\in(1;\infty) \).}
{\( 6 \)}{\( 4 \)}{\( 1 \)}{\( 2 \)}{}{}}
\LANGcs{
\Question{
Určete součet všech kořenů rovnice  \( \bigl| |1-x|-2 \bigr|=1 \) pro \( x\in(1;\infty) \).}
{\( 6 \)}{\( 4 \)}{\( 1 \)}{\( 2 \)}{}{}}
\LANGsk{
\Question{
Určte súčet všetkých koreňov rovnice \( \bigl| |1-x|-2 \bigr|=1 \) pre \( x\in(1;\infty) \).}
{\( 6 \)}{\( 4 \)}{\( 1 \)}{\( 2 \)}{}{}}
\LANGes{
\Question{
Halla la suma de todas las raíces de la ecuación \( \bigl| |1-x|-2 \bigr|=1 \) para \( x\in(1;\infty) \).}
{\( 6 \)}{\( 4 \)}{\( 1 \)}{\( 2 \)}{}{}}
\End

\Data\ID{1003162905}
\Updated{2020-12-29 08:12:03}
\Level{C}
\SubArea{Absolute value equations and inequalities}
\LANGen{
\Question{
Choose the equation that is equivalent to \( \bigl| |2-x|-1 \bigr|=1 \) for \( x\in(3;\infty) \).}
{\( (x-2)-1=1 \)}{\( -(x-2)-1=1 \)}{\( -\bigl(1-(2-x)\bigr)=1 \)}{\( 1-(x-2)=1 \)}{}{}}
\LANGpl{
\Question{
Wybierz równanie równoważne z  \( \bigl| |2-x|-1 \bigr|=1 \) pro \( x\in(3;\infty) \).}
{\( (x-2)-1=1 \)}{\( -(x-2)-1=1 \)}{\( -\bigl(1-(2-x)\bigr)=1 \)}{\( 1-(x-2)=1 \)}{}{}}
\LANGcs{
\Question{
Určete, která z následujících rovnic je ekvivalentní k rovnici \( \bigl| |2-x|-1 \bigr|=1 \) pro \( x\in(3;\infty) \).}
{\( (x-2)-1=1 \)}{\( -(x-2)-1=1 \)}{\( -\bigl(1-(2-x)\bigr)=1 \)}{\( 1-(x-2)=1 \)}{}{}}
\LANGsk{
\Question{
Vyberte rovnicu, ktorá je ekvivalentná s rovnicou \( \bigl| |2-x|-1 \bigr|=1 \) pre \( x\in(3;\infty) \).}
{\( (x-2)-1=1 \)}{\( -(x-2)-1=1 \)}{\( -\bigl(1-(2-x)\bigr)=1 \)}{\( 1-(x-2)=1 \)}{}{}}
\LANGes{
\Question{
Elige la ecuación que es equivalente a \( \bigl| |2-x|-1 \bigr|=1 \) para \( x\in(3;\infty) \).}
{\( (x-2)-1=1 \)}{\( -(x-2)-1=1 \)}{\( -\bigl(1-(2-x)\bigr)=1 \)}{\( 1-(x-2)=1 \)}{}{}}
\End

\Data\ID{1003163002}
\Updated{2022-04-23 05:04:37}
\Level{C}
\SubArea{Absolute value equations and inequalities}
\LANGen{
\Question{
Find all \( t \), \( t\in\mathbb{R} \), such that the following equation with the variable \( x \) has more than two solutions.
\[ \Bigl| |3-x|-3\Bigr|=t \]}
{\( t\in(0;3] \)}{\( t\in\{0;3\} \)}{\( t\in\{3\} \)}{\( t\in[1;3] \)}{}{}}
\LANGpl{
\Question{
Wyznacz wszystkie \( t \), \( t\in\mathbb{R} \), takie dla których podane równanie ze zmienną  \( x \) ma więcej niż dwa rozwiązania.
\[ \Bigl| |3-x|-3\Bigr|=t \]}
{\( t\in(0;3\rangle \)}{\( t\in\{0;3\} \)}{\( t\in\{3\} \)}{\( t\in\langle1;3\rangle \)}{}{}}
\LANGcs{
\Question{
Určete všechna \( t \), \( t\in\mathbb{R} \), pro která má daná rovnice s neznámou \( x \) více než dvě řešení.
\[ \Bigl| |3-x|-3\Bigr|=t \]}
{\( t\in(0;3\rangle \)}{\( t\in\{0;3\} \)}{\( t\in\{3\} \)}{\( t\in\langle1;3\rangle \)}{}{}}
\LANGsk{
\Question{
Určte všetky \( t \), \( t\in\mathbb{R} \), pre ktoré má daná rovnica s neznámou \( x \) viac než dve riešenia.
\[ \Bigl| |3-x|-3\Bigr|=t \]}
{\( t\in(0;3\rangle \)}{\( t\in\{0;3\} \)}{\( t\in\{3\} \)}{\( t\in\langle1;3\rangle \)}{}{}}
\LANGes{
\Question{
Encuentra todos los valores de \( t \), \( t\in\mathbb{R} \), para los que la siguiente ecuación, siendo \( x \) la incognita, tiene más de dos soluciones. 
\[ \Bigl| |3-x|-3\Bigr|=t \]}
{\( t\in(0;3] \)}{\( t\in\{0;3\} \)}{\( t\in\{3\} \)}{\( t\in[1;3] \)}{}{}}
\End

\Data\ID{1003163004}
\Updated{2020-08-23 12:08:54}
\Level{C}
\SubArea{Absolute value equations and inequalities}
\LANGen{
\Question{
Find the sum of all the roots of the given equation.
\[ \Bigl| |x-4|-9\Bigr|=0 \]}
{\( 8 \)}{\( 18 \)}{\( 6 \)}{\( 4 \)}{}{}}
\LANGpl{
\Question{
Wyznacz sumę wszystkich pierwiastków danego równania.
\[ \Bigl| |x-4|-9\Bigr|=0 \]}
{\( 8 \)}{\( 18 \)}{\( 6 \)}{\( 4 \)}{}{}}
\LANGcs{
\Question{
Určete součet všech kořenů dané rovnice.
\[ \Bigl| |x-4|-9\Bigr|=0 \]}
{\( 8 \)}{\( 18 \)}{\( 6 \)}{\( 4 \)}{}{}}
\LANGsk{
\Question{
Určte súčet všetkých koreňov danej rovnice.
\[ \Bigl| |x-4|-9\Bigr|=0 \]}
{\( 8 \)}{\( 18 \)}{\( 6 \)}{\( 4 \)}{}{}}
\LANGes{
\Question{
Halla la suma de todas las raíces de la siguiente ecuación
\[ \Bigl| |x-4|-9\Bigr|=0 \]}
{\( 8 \)}{\( 18 \)}{\( 6 \)}{\( 4 \)}{}{}}
\End

\Data\ID{1003177804}
\Updated{2020-08-23 12:08:28}
\Level{C}
\SubArea{Absolute value equations and inequalities}
\LANGen{
\Question{
Choose the solution set of the given inequality.
\[ \left| |x-5|+2\right| > 3 \]}
{\( (-\infty;4)\cup(6;\infty) \)}{\( (-\infty;-4)\cup(6;\infty) \)}{\( (6;\infty) \)}{\( (-\infty;-5)\cup(6;\infty) \)}{}{}}
\LANGpl{
\Question{
Wybierz zbiór rozwiązań podanej nierówności.
\[ \left| |x-5|+2\right| > 3 \]}
{\( (-\infty;4)\cup(6;\infty) \)}{\( (-\infty;-4)\cup(6;\infty) \)}{\( (6;\infty) \)}{\( (-\infty;-5)\cup(6;\infty) \)}{}{}}
\LANGcs{
\Question{
Určete množinu řešení dané nerovnice.
\[ \left| |x-5|+2\right| > 3 \]}
{\( (-\infty;4)\cup(6;\infty) \)}{\( (-\infty;-4)\cup(6;\infty) \)}{\( (6;\infty) \)}{\( (-\infty;-5)\cup(6;\infty) \)}{}{}}
\LANGsk{
\Question{
Určte množinu riešení danej nerovnice.
\[ \left| |x-5|+2\right| > 3 \]}
{\( (-\infty;4)\cup(6;\infty) \)}{\( (-\infty;-4)\cup(6;\infty) \)}{\( (6;\infty) \)}{\( (-\infty;-5)\cup(6;\infty) \)}{}{}}
\LANGes{
\Question{
Identifica el conjunto de soluciones de la siguiente inecuación.
\[ \left| |x-5|+2\right| > 3 \]}
{\( (-\infty;4)\cup(6;\infty) \)}{\( (-\infty;-4)\cup(6;\infty) \)}{\( (6;\infty) \)}{\( (-\infty;-5)\cup(6;\infty) \)}{}{}}
\End

\Data\ID{1003177805}
\Updated{2020-08-23 01:08:57}
\Level{C}
\SubArea{Absolute value equations and inequalities}
\LANGen{
\Question{
Let \( x\in\mathbb{Z} \). Find the sum of all the solutions of the given inequality.
\[ \Bigl| \bigl| |x+1|+1 \bigr|+1\Bigr| \leq 3 \]}
{\( -3 \)}{\( -2 \)}{\( -4 \)}{\( 0 \)}{}{}}
\LANGpl{
\Question{
Niech \( x\in\mathbb{Z} \). Wyznacz sumę wszystkich rozwiązań podanej nierówności.
\[ \Bigl| \bigl| |x+1|+1 \bigr|+1\Bigr| \leq 3 \]}
{\( -3 \)}{\( -2 \)}{\( -4 \)}{\( 0 \)}{}{}}
\LANGcs{
\Question{
Nechť \( x\in\mathbb{Z} \). Určete součet všech řešení dané nerovnice.
\[ \Bigl| \bigl| |x+1|+1 \bigr|+1\Bigr| \leq 3 \]}
{\( -3 \)}{\( -2 \)}{\( -4 \)}{\( 0 \)}{}{}}
\LANGsk{
\Question{
Nech \( x\in\mathbb{Z} \). Určte súčet všetkých riešení danej nerovnice.
\[ \Bigl| \bigl| |x+1|+1 \bigr|+1\Bigr| \leq 3 \]}
{\( -3 \)}{\( -2 \)}{\( -4 \)}{\( 0 \)}{}{}}
\LANGes{
\Question{
Sea \( x\in\mathbb{Z} \). Halla la suma de todas las raíces de la siguiente inecuación. \[ \Bigl| \bigl| |x+1|+1 \bigr|+1\Bigr| \leq 3 \]}
{\( -3 \)}{\( -2 \)}{\( -4 \)}{\( 0 \)}{}{}}
\End

\Data\ID{1003187302}
\Updated{2020-08-23 07:08:50}
\Level{C}
\SubArea{Absolute value equations and inequalities}
\LANGen{
\Question{
The equation \( \left| |2x-8|-4 \right|=4 \) has:}
{three real solutions}{one real solution}{four real solutions}{two real solutions}{}{}}
\LANGpl{
\Question{
Równanie \( \left| |2x-8|-4 \right|=4 \) ma:}
{trzy rzeczywiste rozwiązania}{jedno rzeczywiste rozwiązanie}{cztery rzeczywiste rozwiązania}{dwa rzeczywiste równania}{}{}}
\LANGcs{
\Question{
Rovnice \( \left| |2x-8|-4 \right|=4 \) má:}
{tři reálná řešení}{jedno reálné řešení}{čtyři reálná řešení}{dvě reálná řešení}{}{}}
\LANGsk{
\Question{
Rovnica \( \left| |2x-8|-4 \right|=4 \) má:}
{tri reálne riešenia}{jedno reálne riešenie}{štyri reálne riešenia}{dve reálne riešenia}{}{}}
\LANGes{
\Question{
La ecuación \( \left| |2x-8|-4 \right|=4 \) tiene:}
{tres soluciones reales}{una solución real}{cuatro soluciones reales}{dos soluciones reales}{}{}}
\End

\Data\ID{1003187401}
\Updated{2022-04-23 05:04:37}
\Level{C}
\SubArea{Absolute value equations and inequalities}
\LANGen{
\Question{
Give the solution set of the equation \( \left| |x-1|-|3-x| \right|=2 \).}
{\( (-\infty;1]\cup[3;+\infty) \)}{\( (-\infty;1)\cup(3;+\infty) \)}{\( \mathbb{R} \)}{\( (1;3) \)}{}{}}
\LANGpl{
\Question{
Wyznacz zbiór rozwiązań równania \( \left| |x-1|-|3-x| \right|=2 \).}
{\( (-\infty;1\rangle\cup\langle3;+\infty) \)}{\( (-\infty;1)\cup(3;+\infty) \)}{\( \mathbb{R} \)}{\( (1;3) \)}{}{}}
\LANGcs{
\Question{
Určete množinu řešení rovnice \( \left| |x-1|-|3-x| \right|=2 \).}
{\( (-\infty;1\rangle\cup\langle3;+\infty) \)}{\( (-\infty;1)\cup(3;+\infty) \)}{\( \mathbb{R} \)}{\( (1;3) \)}{}{}}
\LANGsk{
\Question{
Určte množinu riešení rovnice \( \left| |x-1|-|3-x| \right|=2 \).}
{\( (-\infty;1\rangle\cup\langle3;+\infty) \)}{\( (-\infty;1)\cup(3;+\infty) \)}{\( \mathbb{R} \)}{\( (1;3) \)}{}{}}
\LANGes{
\Question{
Identifica el conjunto de soluciones de la ecuación \( \left| |x-1|-|3-x| \right|=2 \).}
{\( (-\infty;1]\cup[3;+\infty) \)}{\( (-\infty;1)\cup(3;+\infty) \)}{\( \mathbb{R} \)}{\( (1;3) \)}{}{}}
\End

\Data\ID{1003187403}
\Updated{2020-08-23 01:08:51}
\Level{C}
\SubArea{Absolute value equations and inequalities}
\LANGen{
\Question{
How many solutions does the equation \( \left| |2x+6|-8\right|=4 \) have?}
{\( 4 \)}{\( 2 \)}{\( 0 \)}{\( 6 \)}{}{}}
\LANGpl{
\Question{
Ile rozwiązań ma równanie \( \left| |2x+6|-8\right|=4 \)?}
{\( 4 \)}{\( 2 \)}{\( 0 \)}{\( 6 \)}{}{}}
\LANGcs{
\Question{
Kolik řešení má rovnice \( \left| |2x+6|-8\right|=4 \)?}
{\( 4 \)}{\( 2 \)}{\( 0 \)}{\( 6 \)}{}{}}
\LANGsk{
\Question{
Koľko riešení má rovnica \( \left| |2x+6|-8\right|=4 \)?}
{\( 4 \)}{\( 2 \)}{\( 0 \)}{\( 6 \)}{}{}}
\LANGes{
\Question{
¿Cuántas soluciones tiene la ecuación \( \left| |2x+6|-8\right|=4 \)?}
{\( 4 \)}{\( 2 \)}{\( 0 \)}{\( 6 \)}{}{}}
\End

\Data\ID{1003187410}
\Updated{2020-12-29 08:12:14}
\Level{C}
\SubArea{Absolute value equations and inequalities}
\LANGen{
\Question{
How many integers are in the solution set of the given inequalities?
\[ \left|\sqrt{17}-2x\right| \leq 5 \text{ and } \left|2x-\sqrt{17}\right| \leq 5 \]}
{\( 5 \)}{\( 4 \)}{\( 6 \)}{\( 7 \)}{}{}}
\LANGpl{
\Question{
Ile liczb całkowitych należy do zbioru rozwiązań nierówności
\( \left|\sqrt{17}-2x\right| \leq 5 \) i \( \left|2x-\sqrt{17}\right| \leq 5 \)?}
{\( 5 \)}{\( 4 \)}{\( 6 \)}{\( 7 \)}{}{}}
\LANGcs{
\Question{
Kolik celých čísel je v množině řešení daných nerovnic?
\[ \left|\sqrt{17}-2x\right| \leq 5 \text{ a } \left|2x-\sqrt{17}\right| \leq 5 \]}
{\( 5 \)}{\( 4 \)}{\( 6 \)}{\( 7 \)}{}{}}
\LANGsk{
\Question{
Koľko celých čísel je v množine riešení daných nerovníc?
\[ \left|\sqrt{17}-2x\right| \leq 5 \text{ a } \left|2x-\sqrt{17}\right| \leq 5 \]}
{\( 5 \)}{\( 4 \)}{\( 6 \)}{\( 7 \)}{}{}}
\LANGes{
\Question{
¿Cuántos números enteros hay en el conjunto de soluciones de la siguiente inecuación? 
\[ \left|\sqrt{17}-2x\right| \leq 5 \text{ y} \left|2x-\sqrt{17}\right| \leq 5 \]}
{\( 5 \)}{\( 4 \)}{\( 6 \)}{\( 7 \)}{}{}}
\End

\Data\ID{2000002203}
\Updated{2021-12-01 03:12:16}
\Level{C}
\SubArea{Absolute value equations and inequalities}
\LANGen{
\Question{
Identify the solution set of the following inequality.
\[
|x + 2| + |x + 3| \leq 0
\]}
{\( \emptyset \)}{\( \mathbb{R} \)}{\( \mathbb{R} \setminus \{-3;-2\} \)}{\( \{-3;-2\} \)}{}{}}
\LANGpl{
\Question{
Wyznacz zbiór rozwiązań następującej nierówności. 
\[
|x + 2| + |x + 3| \leq 0
\]}
{\( \emptyset \)}{\( \mathbb{R} \)}{\( \mathbb{R} \setminus \{-3;-2\} \)}{\( \{-3;-2\} \)}{}{}}
\LANGcs{
\Question{
Vyberte množinu, která je řešením dané nerovnice.
\[
|x + 2| + |x + 3| \leq 0
\]}
{\( \emptyset \)}{\( \mathbb{R} \)}{\( \mathbb{R} \setminus \{-3;-2\} \)}{\( \{-3;-2\} \)}{}{}}
\LANGsk{
\Question{
Vyberte množinu, ktorá je riešením danej nerovnice.
\[
|x + 2| + |x + 3| \leq 0
\]}
{\( \emptyset \)}{\( \mathbb{R} \)}{\( \mathbb{R} \setminus \{-3;-2\} \)}{\(\{-3;-2\} \)}{}{}}
\LANGes{
\Question{
Identifica el conjunto de soluciones de la siguiente inecuación. 
\[
|x + 2| + |x + 3| \leq 0
\]}
{\( \emptyset \)}{\( \mathbb{R} \)}{\( \mathbb{R} \setminus \{-3;-2\} \)}{\(\{-3;-2\} \)}{}{}}
\End

\Data\ID{2010008502}
\Updated{2022-08-25 10:08:18}
\Level{C}
\SubArea{Absolute value equations and inequalities}
\LANGen{
\Question{
Find all \( t\in\mathbb{R} \), such that the following equation with the variable \( x \) has more than two solutions.
\[ \Bigl| 1-|2-x|\Bigr|+1=t \]}
{\( t\in(1;2 ] \)}{\( t\in\{2\} \)}{\( t\in(2;\infty) \)}{\( t \in \emptyset \)}{}{}}
\LANGpl{
\Question{
Znajdź wszystkie \( t\in\mathbb{R} \), dla których podane równanie ze zmienną \( x \) ma więcej niż dwa rozwiązania.
\[ \Bigl| 1-|2-x|\Bigr|+1=t \]}
{\( t\in(1;2 \rangle \)}{\( t\in\{2\} \)}{\( t\in(2;\infty) \)}{\( t \in \emptyset \)}{}{}}
\LANGcs{
\Question{
Určete všechna  \( t\in\mathbb{R} \), pro která má daná rovnice s neznámou  \( x \) více než dvě řešení.
\[ \Bigl| 1-|2-x|\Bigr|+1=t \]}
{\( t\in(1;2 \rangle \)}{\( t\in\{2\} \)}{\( t\in(2;\infty) \)}{\( t \in \emptyset \)}{}{}}
\LANGsk{
\Question{
Určte všetky  \( t\in\mathbb{R} \), pre ktoré má daná rovnica s neznámou  \( x \) viac než dve riešenia.
\[ \Bigl| 1-|2-x|\Bigr|+1=t \]}
{\( t\in(1;2 \rangle \)}{\( t\in\{2\} \)}{\( t\in(2;\infty) \)}{\( t \in \emptyset \)}{}{}}
\LANGes{
\Question{
Encuentra todos los valores de \( t\in\mathbb{R} \), para los que la siguiente ecuación, siendo \( x \) la incognita, tiene más de dos soluciones.
\[ \Bigl| 1-|2-x|\Bigr|+1=t \]}
{\( t\in(1;2 ] \)}{\( t\in\{2\} \)}{\( t\in(2;\infty) \)}{\( t \in \emptyset \)}{}{}}
\End

\Data\ID{2010008505}
\Updated{2022-08-25 10:08:18}
\Level{C}
\SubArea{Absolute value equations and inequalities}
\LANGen{
\Question{
Give the solution set of the equation \( \left| |2+x|-|x-4| \right|=6\).}
{\( (-\infty;-2] \cup [ 4 ;+\infty) \)}{\( (-\infty;-4] \cup [ 2 ;+\infty) \)}{\( [ -2 ;+\infty) \)}{\(  [ 4 ;+\infty) \)}{}{}}
\LANGpl{
\Question{
Wskaż zbiór rozwiązań równania \( \left| |2+x|-|x-4| \right|=6\).}
{\( (-\infty;-2\rangle \cup \langle 4 ;+\infty) \)}{\( (-\infty;-4\rangle \cup \langle 2 ;+\infty) \)}{\( \langle -2 ;+\infty) \)}{\(  \langle 4 ;+\infty) \)}{}{}}
\LANGcs{
\Question{
Určete množinu řešení rovnice \( \left| |2+x|-|x-4| \right|=6\).}
{\( (-\infty;-2\rangle \cup \langle 4 ;+\infty) \)}{\( (-\infty;-4\rangle \cup \langle 2 ;+\infty) \)}{\( \langle -2 ;+\infty) \)}{\(  \langle 4 ;+\infty) \)}{}{}}
\LANGsk{
\Question{
Určte množinu riešení rovnice \( \left| |2+x|-|x-4| \right|=6\).}
{\( (-\infty;-2\rangle \cup \langle 4 ;+\infty) \)}{\( (-\infty;-4\rangle \cup \langle 2 ;+\infty) \)}{\( \langle -2 ;+\infty) \)}{\(  \langle 4 ;+\infty) \)}{}{}}
\LANGes{
\Question{
Identifica el conjunto de soluciones de la ecuación \( \left| |2+x|-|x-4| \right|=6\).}
{\( (-\infty;-2] \cup [ 4 ;+\infty) \)}{\( (-\infty;-4] \cup [ 2 ;+\infty) \)}{\( [ -2 ;+\infty) \)}{\(  [ 4 ;+\infty) \)}{}{}}
\End

\Data\ID{2010008604}
\Updated{2022-08-25 10:08:18}
\Level{C}
\SubArea{Absolute value equations and inequalities}
\LANGen{
\Question{
Find the sum of all the roots of the given equation.
\[ \bigl|2- |x+1|\bigr|=0 \]}
{\( -2\)}{\( 2\)}{\( -4 \)}{\( 0 \)}{}{}}
\LANGpl{
\Question{
Znajdź sumę wszystkich pierwiastków danego równania.
\[ \bigl|2- |x+1|\bigr|=0 \]}
{\( -2\)}{\( 2\)}{\( -4 \)}{\( 0 \)}{}{}}
\LANGcs{
\Question{
Určete součet všech kořenů dané rovnice.  
\[ \bigl|2- |x+1|\bigr|=0 \]}
{\( -2\)}{\( 2\)}{\( -4 \)}{\( 0 \)}{}{}}
\LANGsk{
\Question{
Určte súčet všetkých koreňov danej rovnice.  
\[ \bigl|2- |x+1|\bigr|=0 \]}
{\( -2\)}{\( 2\)}{\( -4 \)}{\( 0 \)}{}{}}
\LANGes{
\Question{
Halla la suma de todas las raíces de la ecuación.
\[ \bigl|2- |x+1|\bigr|=0 \]}
{\( -2\)}{\( 2\)}{\( -4 \)}{\( 0 \)}{}{}}
\End
